% =============================================================
%  ГОСТ-совместимый шаблон ВКР  (XeLaTeX)
% =============================================================
\documentclass[14pt,a4paper]{extarticle}

% ------- кодировка и язык -------


\usepackage{fontspec}
\usepackage[english,russian]{babel}
\setmainfont[Ligatures=TeX]{Times New Roman}

% ------- геометрия ГОСТ -------
\usepackage[left=3cm,right=1.5cm,top=2cm,bottom=2cm]{geometry}

% ------- полуторный интервал -------
\usepackage{setspace}
\onehalfspacing

% ------- абзацный отступ -------
\usepackage{indentfirst}
\setlength{\parindent}{1.25cm}

% ------- математика -------
\usepackage{amsmath,amssymb}

% ------- таблицы -------
\usepackage{longtable,booktabs,array}
\newcounter{none}
\usepackage{calc}
\usepackage{etoolbox}
\makeatletter
\patchcmd\longtable{\par}{\if@noskipsec\mbox{}\fi\par}{}{}
\makeatother
\IfFileExists{footnotehyper.sty}{\usepackage{footnotehyper}}{\usepackage{footnote}}
\makesavenoteenv{longtable}

% ------- заголовки ГОСТ: по центру -------
\usepackage[center]{titlesec}
\titlelabel{\thetitle.\quad}
% section: прописные, полужирные
\titleformat{\section}{\normalfont\bfseries\centering\Large}{\thesection.}{0.5em}{}
% subsection: строчные, полужирные
\titleformat{\subsection}{\normalfont\bfseries\centering\large}{\thesubsection.}{0.5em}{}

\setcounter{secnumdepth}{3}

% ------- нумерация страниц ГОСТ: внизу по центру -------
\usepackage{fancyhdr}
\pagestyle{fancy}
\fancyhf{}
\renewcommand{\headrulewidth}{0pt}
\cfoot{\thepage}

% ------- нумерация таблиц/рисунков внутри разделов -------
\usepackage{graphicx}
\numberwithin{table}{section}
\numberwithin{figure}{section}
\numberwithin{equation}{section}

% ------- подписи -------
\usepackage[font=singlespacing]{caption}

% ------- гиперссылки -------
\usepackage{bookmark}
\usepackage{hyperref}
\hypersetup{
  pdflang={ru},
  hidelinks,
  pdfcreator={XeLaTeX}}
\IfFileExists{xurl.sty}{\usepackage{xurl}}{}
\urlstyle{same}

% ------- вспомогательное -------
\usepackage{xcolor}
\setlength{\emergencystretch}{3em}
\providecommand{\tightlist}{%
  \setlength{\itemsep}{0pt}\setlength{\parskip}{0pt}}

\author{}
\date{}

\begin{document}

% ===================== ТИТУЛЬНЫЙ ЛИСТ =====================
\thispagestyle{empty}
\begin{center}

МИНИСТЕРСТВО НАУКИ И ВЫСШЕГО ОБРАЗОВАНИЯ РОССИЙСКОЙ ФЕДЕРАЦИИ

\vspace{0.5em}

НОВОСИБИРСКИЙ ГОСУДАРСТВЕННЫЙ УНИВЕРСИТЕТ

\vspace{0.5em}

\textbf{Механико-математический факультет}\\
\textbf{Кафедра вычислительной математики}\\
\textbf{Направление: 02.03.02 «Математика и компьютерные науки»}

\vspace{3cm}

{\Large\textbf{ВЫПУСКНАЯ КВАЛИФИКАЦИОННАЯ РАБОТА}}

\vspace{1cm}

{\large\textbf{Повышение точности регионального прогноза погоды\\
в графовой нейросетевой модели\\
с помощью усвоения данных}}

\vspace{0.5cm}

\textbf{УДК 519.6 : 551.5}

\end{center}

\vspace{2cm}

\noindent
\textbf{Выполнил:} студент 4 курса, группы 22121\\
Табаков Артур Станиславович

\vspace{1em}

\noindent
\textbf{Научный руководитель:}\\
д.ф.-м.н., ВНС ИВМиМГ СО РАН\\
Пененко Алексей Владимирович

\vfill

\begin{center}
\textbf{Новосибирск, 2026}
\end{center}

\newpage

\section*{АННОТАЦИЯ}\label{ux430ux43dux43dux43eux442ux430ux446ux438ux44f}

Графовые нейронные сети стали одним из ведущих инструментов для
краткосрочного и среднесрочного прогноза погоды: такие модели, как
GraphCast, обгоняют по качеству классические численные системы при
значительно меньшей вычислительной стоимости выполнения прогноза. В
настоящей работе исследуется применение архитектуры типа GraphCast
к задаче регионального прогноза погоды на территории Красноярского
края и её интеграция с методами усвоения данных и статистического
постпроцессинга в рамках единого гибридного подхода.

В качестве базовой используется графовая нейросетевая модель формата
«кодировщик--процессор--декодировщик» на икосаэдральной сетке,
воспроизводящая ключевые принципы системы Google GraphCast и обученная
на данных реанализа ERA5. Для регионального прогноза построен
двухуровневый граф: глобальная сетка пониженного разрешения дополняется
детализированной региональной сеткой над интересующей областью, а узлы
двух масштабов связываются общими рёбрами в зоне перекрытия, что
обеспечивает согласованное распространение информации между глобальным
и региональным уровнями и подавляет артефакты на искусственных границах
домена.

При практическом применении модели в авторегрессионном режиме остаются
две ключевые проблемы: накопление ошибки с ростом горизонта прогноза и
систематическое смещение модельных полей относительно наблюдений на
конкретных метеостанциях. Для их решения исследуется гибридная схема,
сочетающая два класса методов. Первый -- усвоение данных на этапе
применения модели: циклическая коррекция прогнозного состояния в
сторону имеющихся наблюдений (nudging) и статистически оптимальное
взвешивание прогноза и наблюдений с учётом пространственных корреляций
ошибок (оптимальная интерполяция). Оба алгоритма применяются поверх
обученной сети и не требуют её переобучения. Второй класс -- статистический
постпроцессинг: на выходе модели обучается регрессионная коррекция по
истории наблюдений на сети метеостанций, исправляющая систематическое
смещение в приземной температуре воздуха.

Проведена серия численных экспериментов на глобальной сетке и для
региона Красноярского края. Выполнено количественное сравнение базовой
GNN-модели, модели с усвоением данных, модели с постпроцессингом и
полной гибридной схемы по стандартным метеорологическим метрикам RMSE,
ACC и bias. Показано, что совместное применение усвоения данных и
постпроцессинга даёт наибольшее снижение накопленной ошибки
авторегрессионного прогноза в ограниченной области, причём вклады этих
двух компонент в значительной степени дополняют друг друга: усвоение
уменьшает рост ошибки по времени, а постпроцессинг -- локальный
систематический bias. Конкретные количественные результаты приведены в
разделах 4--6 настоящей работы.

\textbf{Ключевые слова:} нейросетевой прогноз погоды, графовые нейронные
сети, GraphCast, мультирезолюционный прогноз, усвоение данных,
оптимальная интерполяция, статистический постпроцессинг, MOS, ERA5.


\newpage

{\singlespacing
\tableofcontents
}

\clearpage

\section{ВВЕДЕНИЕ}\label{ux432ux432ux435ux434ux435ux43dux438ux435}

Точный прогноз погоды является одной из наиболее важных и вместе с тем
сложных задач современной науки. От качества прогнозирования зависит
безопасность авиации, принятие сельскохозяйственных решений,
планирование энергетики и управление чрезвычайными ситуациями.
Традиционные методы численного прогноза погоды (Numerical Weather
Prediction, NWP), основанные на численном интегрировании уравнений
гидродинамики атмосферы, достигли высокого уровня точности за несколько
десятилетий развития, однако требуют значительных суперкомпьютерных
ресурсов.

В последние годы произошёл качественный прорыв в применении методов
машинного обучения для прогнозирования погоды. Системы FourCastNet
{[}7{]}, Pangu-Weather {[}8{]}, а в особенности Google GraphCast {[}1{]}
продемонстрировали точность, сопоставимую или превосходящую оперативные
системы ECMWF для ряда ключевых переменных, при несравнимо меньших
вычислительных затратах на этапе прогноза.

Тем не менее большинство существующих нейросетевых систем ориентированы
на глобальный прогноз погоды. Для регионального прогнозирования
возникают специфические трудности, не характерные для глобального
масштаба. Во-первых, при авторегрессионном продлении прогноза ошибки
накапливаются со временем, что особенно критично для ограниченных
областей. Во-вторых, прогноз на ограниченной области не получает
информации о том, что происходит в соседних регионах за пределами
расчётного домена; в результате граничные узлы никогда не «видят»
корректного внешнего окружения, и у краёв области нарастают численные
артефакты, ухудшающие качество прогноза.

Настоящая работа посвящена решению этих двух проблем применительно к
GNN-модели формата GraphCast, адаптированной для регионального прогноза
погоды на территории Красноярского края. Для борьбы с накоплением ошибки
предлагается интегрировать алгоритмы усвоения данных непосредственно в
процедуру прогноза, что позволяет периодически корректировать
прогностическое состояние на основе доступных наблюдений на сети
метеостанций.

\textbf{Цель работы} -- исследовать применимость алгоритмов усвоения
данных (Nudging и оптимальная интерполяция) и статистического
постпроцессинга для повышения точности авторегрессионного
нейросетевого регионального прогноза погоды.

\textbf{Задачи работы:} 1. Реализовать архитектуру GNN-модели прогноза
погоды, инспирированную GraphCast, для регионального домена. 2.
Разработать и реализовать модули Nudging и оптимальной интерполяции для
пост-шагового усвоения данных. 3. Реализовать статистический
постпроцессинг (Learned MOS) для коррекции систематического смещения
приземной температуры по данным сети метеостанций. 4. Провести
численные эксперименты, оценить вклад каждого метода в улучшение метрик
RMSE и ACC. 5. Исследовать зависимость качества усвоения от набора
усваиваемых переменных и гиперпараметров алгоритмов.


\newpage

\section{Обзор литературы}

\subsection{Традиционные методы численного прогноза погоды}\label{ux442ux440ux430ux434ux438ux446ux438ux43eux43dux43dux44bux435-ux43cux435ux442ux43eux434ux44b-ux447ux438ux441ux43bux435ux43dux43dux43eux433ux43e-ux43fux440ux43eux433ux43dux43eux437ux430-ux43fux43eux433ux43eux434ux44b}

Численный прогноз погоды (ЧПП) основан на численном решении системы
уравнений в частных производных, описывающих динамику атмосферы:
уравнений Навье--Стокса, уравнений термодинамики и переноса влаги.
Современные оперативные системы, такие как IFS (ECMWF), GFS (NOAA) и
ICON (DWD), используют спектральные методы (разложение поля по
сферическим гармоникам с точным вычислением производных в спектральном
пространстве) или конечные разности (численная аппроксимация
производных по соседним узлам структурированной сетки).
Пространственное разрешение ведущих глобальных систем составляет
9--25 км, шаг по времени -- несколько минут.

Вместе с численными схемами ЧПП неотделим от системы усвоения данных, задача которой -- совместить модельный прогноз
с реальными наблюдениями. Методы усвоения варьируются от простой
оптимальной интерполяции до более сложных вариационных и ансамблевых
схем, применяемых в оперативных мировых центрах прогноза погоды.

Для регионального масштаба оперативно применяются модели типа WRF,
COSMO, AROME. Они работают в режиме «ограниченной области» (Limited Area
Model, LAM) и требуют задания граничных условий от глобальной модели.
Классический метод нанесения граничных условий -- Davies relaxation
{[}4{]} -- реализует плавное сшивание поля LAM с полем глобальной
модели в буферной зоне.

\subsection{Нейросетевые подходы к прогнозированию погоды}\label{ux43dux435ux439ux440ux43eux441ux435ux442ux435ux432ux44bux435-ux43fux43eux434ux445ux43eux434ux44b-ux43a-ux43fux440ux43eux433ux43dux43eux437ux438ux440ux43eux432ux430ux43dux438ux44e-ux43fux43eux433ux43eux434ux44b}

Применение нейронных сетей для прогноза погоды началось с простых
регрессионных моделей для статистического постпроцессинга {[}9{]}.
Первый прорыв произошёл в 2022--2023 годах, когда возможности
трансформерных архитектур и свёрточных нейросетей были применены
непосредственно для мощного многовременного прогноза.

FourCastNet {[}7{]} -- модель на основе адаптивного
преобразования Фурье (AFNO), обучённая на данных ERA5 с разрешением
0.25°, демонстрирует высокую скорость вычислений при конкурентоспособной
точности краткосрочного прогноза.

Pangu-Weather {[}8{]} -- иерархическая трансформерная модель с
3D Earth Specific Transformer, работающая на разных временных
горизонтах с помощью набора отдельных моделей. Показывает высокую
точность для горизонтов 1--7 суток.

GraphCast {[}1{]} компании Google DeepMind -- GNN-модель на
икосаэдральной сетке, обученная на ERA5 за период 1979--2017. При прогнозе на 10
суток точнее IFS ECMWF по 90\% из 1380 оцениваемых величин. Эта работа
стала ключевым ориентиром для настоящего исследования.

\subsection{GraphCast и архитектуры на основе GNN}\label{graphcast-ux438-ux430ux440ux445ux438ux442ux435ux43aux442ux443ux440ux44b-ux43dux430-ux43eux441ux43dux43eux432ux435-gnn}

Графовые нейронные сети позволяют работать с нерегулярными
сетками -- то есть сетками, не имеющими фиксированной структуры
строк/столбцов, где соседи узла задаются произвольным графом, а не
индексами вида \([i\pm 1, j\pm 1]\). Они естественным образом
моделируют передачу информации между пространственно связанными
узлами. В контексте прогноза погоды GNN привлекательны тем, что:
\begin{itemize}
\item позволяют использовать глобальную икосаэдральную мозаику,
которая почти равномерно покрывает сферу, в отличие от географической
сетки с неравномерным шагом по широте;
\item обеспечивают масштабируемость через иерархию разрешений;
\item могут эффективно обрабатывать пространственно ограниченные
области (региональный прогноз).
\end{itemize}

Архитектура GraphCast реализует трёхуровневую схему: кодировщик
(grid в mesh), процессор (взаимодействие узлов mesh) и
декодировщик (mesh в grid). Процессор реализован как стек «Interaction Networks»
{[}10{]}: на каждом шаге обновляются сначала признаки рёбер, затем --
узлов, причём каждый шаг использует \emph{residual-соединение}
-- арифметическое сложение входа блока с выходом обучаемого
преобразования по схеме \(y = x + F(x)\). Такая конструкция (введённая
в ResNet) ускоряет сходимость глубоких сетей и позволяет каждому шагу
Processor учить лишь поправку к латентному состоянию, а не
перевычислять его целиком. Подход имеет теоретическое обоснование как
аппроксимация численного интегрирования дифференциальных уравнений
через message passing {[}11{]}.

\subsection{Методы усвоения данных в прогнозировании погоды}\label{ux43cux435ux442ux43eux434ux44b-ux443ux441ux432ux43eux435ux43dux438ux44f-ux434ux430ux43dux43dux44bux445-ux432-ux43fux440ux43eux433ux43dux43eux437ux438ux440ux43eux432ux430ux43dux438ux438-ux43fux43eux433ux43eux434ux44b}

Nudging (метод подталкивания) -- один из наиболее ранних и
практически применяемых методов усвоения данных в мезомасштабных моделях
{[}5{]}. Основная идея состоит в добавлении к правой части уравнения
прогноза члена, пропорционального разности между модельным полем и
наблюдением:
\[\frac{\partial x}{\partial t} = F(x) + \alpha \cdot (x^{obs} - x)\]
где \(\alpha\) -- коэффициент подталкивания, \(x^{obs}\) --
наблюдаемое значение. В дискретной форме, применяемой в данной работе на
этапе выдачи прогноза:
\[x_{t+1}^{(new)} = x_{t+1} + \alpha (x_{t+1}^{obs} - x_{t+1})\] Метод
прост в реализации, требует лишь одного гиперпараметра (\(\alpha\)), но
применяет правку только в точках наблюдений -- пространственная
корреляция ошибок не учитывается.

Оптимальная интерполяция (OI) {[}2{]} -- классический метод
статистического оценивания, минимизирующий дисперсию ошибки анализа.
Алгоритм получает взвешенное сочетание прогноза и наблюдений с учётом
пространственных корреляций ошибок: \[x_a = x_b + K(y_o - Hx_b)\]
\[K = BH^T(HBH^T + R)^{-1}\] где \(x_b\) -- фоновое поле (прогноз),
\(x_a\) -- анализ, \(y_o\) -- наблюдения, \(H\) -- оператор
наблюдений, \(B\) -- матрица ковариации ошибок прогноза, \(R\) --
матрица ковариации ошибок наблюдений, \(K\) -- матрица усиления
Калмана. OI обеспечивает пространственное «размазывание» информации от
редких станций на соседние узлы, что делает её значительно более мощным
инструментом в сравнении с Nudging при разреженной сети наблюдений.


\newpage

\section{ОПИСАНИЕ МОДЕЛИ}\label{ux43eux43fux438ux441ux430ux43dux438ux435-ux43cux43eux434ux435ux43bux438}

\subsection{Постановка задачи и входные данные}\label{ux43fux43eux441ux442ux430ux43dux43eux432ux43aux430-ux437ux430ux434ux430ux447ux438-ux438-ux432ux445ux43eux434ux43dux44bux435-ux434ux430ux43dux43dux44bux435}

Задача прогноза формулируется как авторегрессионное отображение: из
\(T_{obs}\) последовательных временных срезов состояния атмосферы
предсказывается состояние на \(T_{pred}\) шагов вперёд. Каждый временной
срез задаётся полем из \(C\) метеорологических переменных на регулярной
координатной сетке.

В качестве обучающих данных использован реанализ ERA5 {[}6{]} от ECMWF
за период 2010--2019 годов с временным шагом 6 часов. Для глобальной
модели используется сетка 512$\times$256 узлов, для регионального
дообучения~-- более мелкая сетка с шагом 0.25° над выделенной
областью (см. раздел~7). Здесь \(T_{obs}\) обозначает длину окна
наблюдений (число последовательных временных срезов, подаваемых на
вход модели), а \(T_{pred}\) -- длину окна прогноза (число шагов,
на которые предсказывается состояние). Набор из 19 динамических и
статических переменных включает:

{\def\LTcaptype{none} % do not increment counter
\begin{longtable}[]{@{}
  >{\raggedright\arraybackslash}p{(\linewidth - 2\tabcolsep) * \real{0.5000}}
  >{\raggedright\arraybackslash}p{(\linewidth - 2\tabcolsep) * \real{0.5000}}@{}}
\toprule\noalign{}
\begin{minipage}[b]{\linewidth}\raggedright
Переменная
\end{minipage} & \begin{minipage}[b]{\linewidth}\raggedright
Описание
\end{minipage} \\
\midrule\noalign{}
\endhead
\bottomrule\noalign{}
\endlastfoot
\emph{t2m} & Температура воздуха на высоте 2 м \\
\emph{msl} & Давление на уровне моря \\
\emph{10u}, \emph{10v} & Компоненты ветра на высоте 10 м \\
\emph{tp} & Суммарные осадки за 6 ч \\
\emph{sp} & Поверхностное давление \\
\emph{tcwv} & Общее содержание водяного пара в столбе атмосферы \\
\emph{z\_surf} & Геопотенциал поверхности (рельеф, статика) \\
\emph{lsm} & Маска суши/моря (статика) \\
\emph{t}850, \emph{u}850, \emph{v}850, \emph{z}850, \emph{q}850 &
Температура, компоненты ветра, геопотенциал и удельная влажность на
уровне 850 гПа \\
\emph{t}500, \emph{u}500, \emph{v}500, \emph{z}500, \emph{q}500 &
То же на уровне 500 гПа \\
\end{longtable}
}

Кроме динамических переменных, каждый узел сетки снабжается
статическими признаками: нормализованными значениями широты
и долготы (приведёнными к диапазону \([-1, 1]\) делением на 90° и 180°
соответственно), а также их синус/косинус-преобразованиями --
стандартным способом кодирования периодических величин, при котором
одна угловая координата \(\alpha\) заменяется парой \((\sin\alpha,
\cos\alpha)\), что позволяет модели правильно учитывать близость
значений вблизи разрыва \(0° \leftrightarrow 360°\). Полученный набор
статических признаков объединяется с динамическими признаками узла
перед подачей в кодировщик и не меняется при прогнозировании.

Конфигурация обучения: \(T_{obs} = 4\) временных шага (24 ч истории),
\(T_{pred} = 4\) (прогноз +6, +12, +18, +24 ч).

\subsection{Икосаэдральная мозаичная сетка}\label{ux438ux43aux43eux441ux430ux44dux434ux440ux430ux43bux44cux43dux430ux44f-ux43cux43eux437ux430ux438ux447ux43dux430ux44f-ux441ux435ux442ux43aux430}

Для создания латентного пространства модели используется икосаэдральная
мозаичная сетка (icosahedral mesh), обеспечивающая квазиравномерное
покрытие сферы. Исходный икосаэдр (20 равносторонних треугольников, 12
вершин) последовательно детализируется: на каждом уровне каждый
треугольник делится на 4 подтреугольника путём вставки вершин в середины
рёбер с последующей нормировкой на сферу. Уровень детализации \(L\)
даёт \(10 \cdot 4^L + 2\) вершин.

В настоящей работе используется иерархия двух уровней сетки: \(L = 3\)
(642 вершины) и \(L = 5\) (10 242 вершины). В процессоре модель
оперирует на самом мелком уровне \(L = 5\). Особенности обрезки
мозаики до регионального домена обсуждаются в разделе~7.

\subsection{Построение графов соединений}\label{ux43fux43eux441ux442ux440ux43eux435ux43dux438ux435-ux433ux440ux430ux444ux43eux432-ux441ux43eux435ux434ux438ux43dux435ux43dux438ux439}

Модель использует три графа:

Граф кодировщика. Каждый узел регулярной сетки
соединяется с ближайшими mesh-узлами по радиусному критерию: радиус
связи равен \(r = q \cdot \ell_{max}\), где \(\ell_{max}\) -- длина
самого длинного ребра мозаики на тонком уровне, \(q = 0.6\) (параметр
конфигурации). Это обеспечивает <<закачку>> информации из плотной
географической сетки в разреженную мозаику.

Граф процессора. Граф строится по рёбрам мозаики
\(L = 5\). Каждое ребро снабжается пространственными признаками
-- 4-мерным вектором: расстояние между концами и их 3D-координаты в
декартовой системе отсчёта. Эти признаки передаются в InteractionNet при
обновлении рёбер.

Граф декодировщика. Каждый узел географической
сетки проецируется в треугольник мозаики, в котором он лежит, и получает
входные сигналы от трёх вершин этого треугольника.

\subsection{Архитектура нейронной сети}\label{ux430ux440ux445ux438ux442ux435ux43aux442ux443ux440ux430-ux43dux435ux439ux440ux43eux43dux43dux43eux439-ux441ux435ux442ux438}

Основной блок модели реализует конвейер кодировщик--процессор--декодировщик.

Кодировщик. Принимает конкатенацию динамических и
статических признаков для всех \((N_{grid} + N_{mesh})\) узлов. Состоит
из полносвязного перцептрона (два скрытых слоя по 128 нейронов,
активация Swish\footnote{Swish $\sigma(x)\cdot x$, где $\sigma$~--
сигмоида; гладкая нелинейность, ведущая себя как ReLU при больших
положительных значениях и плавно затухающая при отрицательных.},
послойная нормализация LayerNorm\footnote{LayerNorm~-- послойная
нормализация: для каждого узла признаки центрируются и масштабируются
по их собственному среднему и дисперсии вдоль канала признаков, после
чего применяются обучаемые масштаб и сдвиг. В отличие от BatchNorm не
зависит от размера батча и устойчиво работает при графовом обучении.})
и слоя графовой свёртки GCN (два слоя по 128 нейронов). Выход --
латентный вектор размерности 128 для каждого узла.

Процессор. Работает только на узлах мозаичной сетки
\(N_{mesh}\). Представляет собой стек из 8 последовательных шагов
Interaction Network с неразделёнными весами (т.\,е. на каждом шаге
обучается собственный набор параметров MLP). Каждый шаг состоит из
следующих операций:
\begin{enumerate}
\def\labelenumi{\arabic{enumi}.}
\item Обновление рёбер:
\(e_{ij}' = \text{MLP}_{edge}([h_i \| h_j \| e_{ij}]) + e_{ij}\) --
конкатенация признаков отправителя, получателя и текущего ребра, затем
MLP и остаточное соединение (residual connection).
\item Агрегация: для каждого узла-получателя усредняются
обновлённые признаки всех входящих рёбер.
\item Обновление узлов:
\(h_i' = \text{MLP}_{node}([h_i \| \overline{e}_i]) + h_i\) --
конкатенация с агрегированными признаками рёбер, MLP и остаточное
соединение.
\item LayerNorm после каждого шага.
\end{enumerate}
Размер скрытого состояния узлов и рёбер в латентном пространстве --
128. Исходные признаки рёбер четырёхмерны (евклидово расстояние
между узлами и три компоненты единичного вектора направления в
3D-декартовых координатах); они проецируются обучаемым линейным слоём
в пространство размерности 128 перед первым шагом message passing.

Декодировщик. Симметрично кодировщику, выполняет
обратное message passing с помощью лёгкого слоя графовой свёртки GCN
(тяжёлое обновление рёбер через InteractionNetwork сосредоточено в
процессоре). Получает конкатенацию grid-латентов (от кодировщика) и
mesh-латентов (от процессора); для каждого узла регулярной сетки
агрегируются признаки от вершин ближайшего треугольника мозаики.
Финальный линейный слой проецирует в \(C \cdot T_{pred}\) выходных
каналов -- предсказания на все горизонты сразу.

Итоговый граф вычислений описывается следующей последовательностью
преобразований. Введём обозначения: \(G = N_{grid}\)~-- число узлов
регулярной географической сетки, \(M = N_{mesh}\)~-- число узлов
мозаичной сетки на тонком уровне, \(C\)~-- число метеорологических
переменных в одном временном срезе, \(C_s\)~-- число статических
признаков на узел, \(d = 128\)~-- размерность латентного пространства,
\(\mathbf{X}^{(0)}\)~-- матрица входных признаков всех узлов,
\(\mathbf{Z}\)~-- латентное представление узлов после кодировщика,
\(\hat{\mathbf{X}}\)~-- выход модели (поля прогноза на регулярной
сетке). Тогда:

\[\mathbf{X}^{(0)} \in \mathbb{R}^{(G+M)\times(T_{obs}\cdot C + C_s)} \xrightarrow{\text{Encoder}} \mathbf{Z} \in \mathbb{R}^{(G+M)\times d}\]

\[\mathbf{Z}_{mesh} \in \mathbb{R}^{M\times d} \xrightarrow{\text{Processor}\times 8} \tilde{\mathbf{Z}}_{mesh} \in \mathbb{R}^{M\times d}\]

\[[\mathbf{Z}_{grid} \,\|\, \tilde{\mathbf{Z}}_{mesh}] \in \mathbb{R}^{(G+M)\times d} \xrightarrow{\text{Decoder}} \hat{\mathbf{X}} \in \mathbb{R}^{G\times (C\cdot T_{pred})}\]

где \(\|\) обозначает конкатенацию матриц по оси узлов, а из выхода
декодировщика берутся только первые \(G\) строк, соответствующие узлам
регулярной сетки.

\subsection{Функция потерь и обучение}\label{ux444ux443ux43dux43aux446ux438ux44f-ux43fux43eux442ux435ux440ux44c-ux438-ux43eux431ux443ux447ux435ux43dux438ux435}

Модель обучается методом минимизации взвешенного MSE:

\[\mathcal{L} = \frac{1}{N \cdot C \cdot T} \sum_{n,g,c,t} w_g \, m_g \, \left(\hat{x}_{n,g,c,t} - x_{n,g,c,t}\right)^2\]

где \(N\)~-- число обучающих примеров на одной
итерации, \(C\)~-- число прогнозируемых переменных, \(T = T_{pred}\)~--
число шагов прогноза, индексы \(n, g, c, t\) пробегают соответственно
примеры батча, узлы регулярной сетки, переменные и горизонты прогноза;
\(w_g = \cos(\varphi_g) / \overline{\cos(\varphi)}\) -- широтные весовые
коэффициенты, выравнивающие вклад полюсов и экватора, \(m_g\)~--
пространственная маска, обнуляющая буферную зону.

Для борьбы с накоплением ошибки применяется поэтапное наращивание
горизонта прогноза: в первых эпохах модель предсказывает 1 шаг, затем
длина авторегрессионной развёртки постепенно увеличивается до
\(T_{pred} = 4\). Оптимизатор -- Adam с начальной скоростью обучения
\(3 \cdot 10^{-4}\), ранняя остановка по метрике на валидационной
выборке.


\newpage

\section{УСВОЕНИЕ ДАННЫХ НА ЭТАПЕ ПРОГНОЗА}\label{sec:da-inference}

\subsection{Постановка задачи усвоения}\label{ux43fux43eux441ux442ux430ux43dux43eux432ux43aux430-ux437ux430ux434ux430ux447ux438-ux443ux441ux432ux43eux435ux43dux438ux44f}

При авторегрессионном применении региональной модели ошибки нарастают с
каждым прогностическим шагом. Если в момент времени \(t_k\) применить
доступные наблюдения для коррекции модельного поля, накопленная ошибка
будет снижена перед следующим шагом интегрирования. Формально задача
усвоения состоит в нахождении анализа \(x_a\), оптимально
объединяющего прогностическое фоновое поле \(x_b\) (выход
модели) и наблюдения \(y_o\).

В рамках настоящей работы усвоение выполняется между
последовательными прогностическими шагами модели. Это соответствует схеме
оперативного цикла усвоения при непрерывном прогнозе:
предсказание, коррекция наблюдениями, подача скорректированного
состояния на следующий шаг.

Для имитации сети наблюдений использовалась субдискретизация поля ERA5: случайным образом отбирается
10\% узлов сетки; именно в этих условных станциях прогнозное
поле сравнивается с реальными значениями ERA5. Вне узлов станций поле
считается ненаблюдаемым (значение отсутствует). Были сформированы
четыре набора наблюдений, отличающихся набором включённых
переменных:

{\def\LTcaptype{none} % do not increment counter
\begin{longtable}[]{@{}
  >{\raggedright\arraybackslash}p{(\linewidth - 4\tabcolsep) * \real{0.3333}}
  >{\raggedright\arraybackslash}p{(\linewidth - 4\tabcolsep) * \real{0.3333}}
  >{\raggedright\arraybackslash}p{(\linewidth - 4\tabcolsep) * \real{0.3333}}@{}}
\toprule\noalign{}
\begin{minipage}[b]{\linewidth}\raggedright
Набор
\end{minipage} & \begin{minipage}[b]{\linewidth}\raggedright
Переменные
\end{minipage} & \begin{minipage}[b]{\linewidth}\raggedright
Физическая интерпретация
\end{minipage} \\
\midrule\noalign{}
\endhead
\bottomrule\noalign{}
\endlastfoot
Все переменные & все 19 & Эталонный сценарий \\
Приземные & \emph{t2m, msl, 10u, 10v, tp} & Простая метеостанция \\
Динамические & \emph{10u, 10v, u, v, z} (все уровни), \emph{msl} &
Динамические переменные \\
Температура & \emph{t2m, t}850, \emph{t}500 & Только температура \\
\end{longtable}
}

\subsection{Метод подталкивания (Nudging)}\label{ux43cux435ux442ux43eux434-ux43fux43eux434ux442ux430ux43bux43aux438ux432ux430ux43dux438ux44f-nudging}

Алгоритм. Пусть \(x_t \in \mathbb{R}^{G \times C}\) --
прогностическое поле (G узлов, C переменных) после \(t\)-го шага модели,
\(y_t \in \mathbb{R}^{G \times C}\) -- поле наблюдений с пропусками
в ненаблюдаемых узлах. Тогда анализ вычисляется по формуле:

\[x_a(g, c) = \begin{cases}
x_b(g,c) + \alpha \bigl(y_o(g,c) - x_b(g,c)\bigr), & \text{если наблюдение доступно}, \\
x_b(g,c), & \text{иначе}.
\end{cases}\]

Параметр \(\alpha \in [0, 1]\) управляет силой усвоения: \(\alpha = 0\)
-- прогноз без коррекции, \(\alpha = 1\) -- полная замена прогноза
наблюдением.

Режим «Sequential Nudging». В данной работе применяется
последовательная коррекция: после каждого авторегрессионного шага
модельный выход корректируется на основе
соответствующего среза наблюдений, после чего скорректированное
состояние подаётся на следующий шаг.

Схема алгоритма. Процедура Sequential Nudging формально
описывается следующим образом. Пусть
\(\mathcal{O} \subset \{1,\ldots,G\}\) -- множество узлов наблюдений
(ненулевых в \(y_o\)).

\[x_a(g, c) = \begin{cases} x_b(g,c) + \alpha\,\delta(g,c), & g \in \mathcal{O},\\ x_b(g,c), & g \notin \mathcal{O}, \end{cases} \quad \delta(g,c) = y_o(g,c) - x_b(g,c).\]

Выбор \(\alpha\). В ходе экспериментов по подбору
гиперпараметров было проверено \(\alpha \in \{0.1, 0.5, 0.9, 1.3\}\).
Оптимальное значение \(\alpha = 0.5\) показало наилучший баланс между
коррекцией ошибки и устойчивостью -- см. раздел 6.5.

\subsection{Оптимальная интерполяция}\label{ux43eux43fux442ux438ux43cux430ux43bux44cux43dux430ux44f-ux438ux43dux442ux435ux440ux43fux43eux43bux44fux446ux438ux44f}

Теория. Оптимальная интерполяция является частным случаем
метода BLUE (Best Linear Unbiased Estimator) {[}4{]}, который при
гауссовских предположениях даёт оценку с минимальной дисперсией
ошибки анализа среди всех линейных несмещённых оценок. Матрица
усиления Калмана \(K\) выводится аналитически из условия минимума
следа ковариационной матрицы анализа:

\[K = B H^T (H B H^T + R)^{-1}\]

Поправка анализа к фону (величина \(y_o - H x_b\) называется
инновацией -- разницей между реальным наблюдением и его
прогнозным аналогом):
\[x_a = x_b + K (y_o - H x_b)\]

Матрица ковариации ошибок фона \(B\). Пространственные
корреляции ошибок прогноза задаются гауссовской моделью:
\[B_{ij} = \sigma_b^2 \cdot \exp\left(-\frac{d_{ij}^2}{L^2}\right)\] где
\(d_{ij}\) -- расстояние по поверхности Земли между узлами \(i\) и
\(j\) (в метрах по формуле гаверсинуса), \(L\) -- радиус корреляции
(гиперпараметр), \(\sigma_b\) -- стандартное отклонение ошибки фона.

Физический смысл. Если в точке \(A\) прогноз отклонился от
наблюдения на \(\delta\), то ОИ «размазывает» эту поправку на соседние
узлы с весом, убывающим с увеличением расстояния. Пространственный
масштаб убывания определяется параметром \(L\).

Матрица ошибок наблюдений \(R\). Принята диагональной:
\(R = \sigma_o^2 \cdot I_{N_{obs}}\), что соответствует предположению о
независимости ошибок разных наблюдений.

Оператор наблюдений \(H\). Матрица
\(H \in \{0,1\}^{N_{obs} \times G}\) -- билинейный оператор проекции:
для каждого наблюдения находится ближайший узел сетки и соответствующий
элемент \(H\) устанавливается в 1.

Вычислительная схема. Матрица \(B \in \mathbb{R}^{G \times G}\)
строится единожды при инициализации алгоритма. Для каждой переменной
\(c = 1,\ldots,C\) на каждом шаге коррекции выполняется:

\[K^{(c)} = B H^{(c)\,T}\bigl(H^{(c)} B H^{(c)\,T} + R^{(c)} + \varepsilon I\bigr)^{-1}\]

\[x_a^{(c)} = x_b^{(c)} + K^{(c)}\bigl(y_o^{(c)} - H^{(c)} x_b^{(c)}\bigr)\]

где \(\varepsilon = 10^{-5}\) -- параметр регуляризации, гарантирующий
невырожденность обращаемой матрицы. Оператор наблюдений \(H^{(c)}\)
строится для каждой переменной независимо: ненулевые строки
соответствуют наблюдаемым узлам, матрица
\(R^{(c)} = \sigma_o^2 I_{N_{obs}}\).

Параметры. Ключевые гиперпараметры ОИ: \(\sigma_b\) (фиксирован
= 0.8), \(\sigma_o\) (подбирался из \(\{0.2, 0.5, 0.8\}\)), \(L\)
(подбирался из \(\{50, 150, 300\}\) км для пилотных экспериментов
и далее из расширенного диапазона \(\{10, 25, 50, 100, 150, 200, 300, 500\}\) км
для мультирезолюционной модели). Матрица \(B\) одинакова
для всех переменных, поэтому коррекция применяется независимо по каждому
каналу \(c\).


\newpage

\section{ОБРАБОТКА ГРАНИЦ РАСЧЁТНОЙ ОБЛАСТИ}\label{ux43eux431ux440ux430ux431ux43eux442ux43aux430-ux433ux440ux430ux43dux438ux446-ux440ux430ux441ux447ux451ux442ux43dux43eux439-ux43eux431ux43bux430ux441ux442ux438}

\subsection{Проблема искусственных границ}\label{ux43fux440ux43eux431ux43bux435ux43cux430-ux438ux441ux43aux443ux441ux441ux442ux432ux435ux43dux43dux44bux445-ux433ux440ux430ux43dux438ux446}

Данный раздел описывает первый из двух рассматриваемых в работе
подходов к подавлению граничных артефактов в региональной GNN-модели
-- явное сшивание вырезанной региональной области с внешним полем в
буферной зоне по схеме Davies, апплицируемое после получения
прогноза. Альтернативный подход -- включение внешнего глобального
поля непосредственно в входные данные GNN в виде единой
мультирезолюционной сетки -- рассматривается в разделе~7. Пилотные
результаты метода Davies из этого раздела используются в пилотной серии
экспериментов в разделе~6.

При региональном прогнозировании граница расчётного домена является
искусственной: на ней отсутствуют физические условия, а
информационный поток из внешней области обрывается. Это приводит к
нарастанию ошибок у границ и отражению волн, которые в реальности должны
выходить из домена (проблема «radiation condition»).

В рамках нейросетевого подхода проблема границ проявляется ещё резче:
модель обучена на глобальных или иных данных, но применяется к
вырезанной области -- и граничные узлы никогда не «видят» правильного
внешнего окружения.

Классическое решение в задачах NWP -- метод Davies {[}4{]}:
в узкой буферной зоне на границе ROI прогноз региональной модели
\(x_{LAM}\) плавно сшивается с фоновым полем глобальной модели
\(x_{bg}\) при помощи весовой функции, монотонно убывающей от центра
к границе. Эта схема подавляет нефизичные отражения волн от стенок
домена и согласует решение с условиями вне ROI:

\[x = w \cdot x_{LAM} + (1 - w) \cdot x_{bg}\]

где весовая функция \(w\) монотонно растёт от 0 у внешней границы до 1 в
центре домена.

\subsection{Буферная зона с косинусным сглаживанием}\label{ux431ux443ux444ux435ux440ux43dux430ux44f-ux437ux43eux43dux430-ux441-ux43aux43eux441ux438ux43dux443ux441ux43dux44bux43c-ux441ux433ux43bux430ux436ux438ux432ux430ux43dux438ux435ux43c}

В данной работе реализовано двумерное окно Ханна:

\[w_{lon}(i) = \begin{cases} \frac{1}{2}\left[1 - \cos\!\left(\frac{\pi i}{b}\right)\right], & 0 \le i < b, \\ 1, & b \le i < N_{lon} - b, \\ \frac{1}{2}\left[1 - \cos\!\left(\frac{\pi (N_{lon}-1-i)}{b}\right)\right], & N_{lon} - b \le i < N_{lon} \end{cases}\]

и аналогично \(w_{lat}(j)\). Итоговая двумерная маска:
\[W(i, j) = w_{lon}(i) \cdot w_{lat}(j)\]

Выбор косинусного окна вместо линейного обусловлен следующими
соображениями:

\begin{enumerate}
\def\labelenumi{\arabic{enumi}.}
\item Производная \(W\) обращается в нуль на обоих концах буферной
зоны -- обеспечивается \(C^1\)-гладкость перехода, что важно для
стабильности нейросетевого предсказания на следующем шаге.
\item Косинусное окно подавляет спектральные утечки --
перераспределение энергии резкого перехода в
высокочастотные гармоники, которые модель воспринимает
как ложные мелкомасштабные структуры.
\item Косинусное окно соответствует принципу максимально плавного
сопряжения полей, изложенному в работе Davies {[}4{]}.
\end{enumerate}

Процедура применения. На каждом шаге авторегрессионного
прогноза:
\[x_{out}(g) = W(g) \cdot x_{LAM}(g) + (1 - W(g)) \cdot x_{bg}(g)\]

где \(x_{bg}\) -- фоновое поле (в экспериментах -- данные ERA5, в
оперативном режиме -- прогноз глобальной модели). Значение маски
\(W(g) \in [0, 1]\) плоского массива узлов вычисляется по заранее
построенному двумерному окну Ханна с шириной буфера \(b = 5\) узлов.


\newpage

\section{ВЫЧИСЛИТЕЛЬНЫЕ ЭКСПЕРИМЕНТЫ (ПИЛОТНОЕ
ИССЛЕДОВАНИЕ)}\label{ux432ux44bux447ux438ux441ux43bux438ux442ux435ux43bux44cux43dux44bux435-ux44dux43aux441ux43fux435ux440ux438ux43cux435ux43dux442ux44b}

Данная глава описывает пилотную серию экспериментов по усвоению данных,
выполненную на ранней автономной региональной модели над Западной
Сибирью с прицельной оценкой качества над Новосибирской областью
(НСКО). Параметры конфигурации сведены в таблицу~3 ниже. Эта серия
позволила отработать алгоритмы Nudging и ОИ, выбрать диапазоны
гиперпараметров и оценить относительную эффективность подходов; её
результаты служат отправной точкой для основной серии экспериментов на
мультирезолюционной модели над Красноярском (раздел~7.6) и не должны
напрямую сопоставляться с числами из последующих глав, поскольку
измерены на иной территории, более грубой сетке и другом тестовом
подмножестве.

Кратко о схеме работы пилотной модели. \emph{Обучающая сетка}~--
глобальная 64$\times$32 узла, покрывающая весь земной шар (шаг около
5.6°); веса GNN-модели подбираются именно на ней. \emph{Полный домен}~--
прямоугольный региональный фрагмент над Западной Сибирью
(47--57.5°\,с.\,ш., 75--90°\,в.\,д., сетка 61$\times$41 узел, шаг 0.25°),
вырезанный из ERA5 высокого разрешения 1440$\times$721; на нём
выполняется прогноз уже обученной модели и считаются основные метрики.
\emph{НСКО}~-- сокращение ``Новосибирская область'', центральный
подрегион полного домена (53--57°\,с.\,ш., 79--85°\,в.\,д., 425 узлов),
вокруг которого расположены условные станции; на этом подрегионе
считаются прицельные метрики, поскольку именно он представляет собой
целевую территорию применения модели. Таким образом, обучение и оценка
ведутся на разных сетках: веса, выученные на глобальной 64$\times$32,
применяются к региональному фрагменту 61$\times$41 без дообучения.

\subsection{Данные и конфигурация экспериментов}\label{ux434ux430ux43dux43dux44bux435-ux438-ux43aux43eux43dux444ux438ux433ux443ux440ux430ux446ux438ux44f-ux44dux43aux441ux43fux435ux440ux438ux43cux435ux43dux442ux43eux432}

Модель обучалась на данных ERA5 за период 2010--2019 гг. (тренировочная
выборка), 2020 первое полугодие (валидация), второе полугодие 2020 года
(тестирование). Конфигурация использованной модели:

{\def\LTcaptype{none} % do not increment counter
\begin{longtable}[]{@{}ll@{}}
\toprule\noalign{}
Параметр & Значение \\
\midrule\noalign{}
\endhead
\bottomrule\noalign{}
\endlastfoot
Обучающая сетка & глобальная 64$\times$32 узла (весь земной шар) \\
Сетка инференса и оценки (полный домен) & 61$\times$41 узел, шаг 0.25° \\
География полного домена & 47--57.5°\,с.\,ш., 75--90°\,в.\,д. (Западная Сибирь) \\
Целевой подрегион (НСКО) & 53--57°\,с.\,ш., 79--85°\,в.\,д. (425 узлов) \\
Источник исходных полей & ERA5 1440$\times$721 (региональный вырез) \\
Обрезка мозаики (буфер) & 15° вокруг полного домена \\
Переменных & 15 динамических \\
\(T_{obs}\) & 4 шага (24 ч предыстории) \\
\(T_{pred}\) & 4 шага (+6, +12, +18, +24 ч) \\
Уровни мозаики & L=3, L=5 \\
Размерность латента & 128 \\
Шагов процессора & 8 \\
Эпох обучения & 50 \\
Оптимизатор & Adam, lr=\(3\cdot10^{-4}\) \\
\end{longtable}
}

Икосаэдральная мозаика для регионального графа обрезается до полного
домена с запасом в 15° по каждой стороне: дополнительная буферная
полоса вне расчётной сетки сохраняется в графе процессора, чтобы при
\(N_{mp}=8\) шагах message passing внутренние узлы получали корректный
информационный приток от соседей и не страдали от «эффекта стены» на
границе. По сравнению с глобальным графом число mesh-узлов и рёбер
сокращается на порядок, что и делает региональное обучение
выполнимым на одной GPU.

Тестовая выборка содержала 13 независимых примеров. Для оценки
усвоения данных использовалась сеть из 250 условных станций
(10\% узлов сетки), выбранных случайно, но фиксированных для всех
экспериментов. В качестве «наблюдений» брались значения ERA5 в этих
узлах: это идеализированный сценарий, где данные имитируются по
реанализу, что позволяет сравнивать методы усвоения данных в контролируемых
условиях.

\subsection{Метрики качества}\label{ux43cux435ux442ux440ux438ux43aux438-ux43aux430ux447ux435ux441ux442ux432ux430}

Для оценки прогнозов использовались стандартные метеорологические
метрики:

RMSE (Root Mean Squared Error) -- в нормированных единицах
(после стандартизации по каналу):
\[\text{RMSE} = \sqrt{\frac{1}{G \cdot C} \sum_{g,c} \left(\hat{x}_{g,c} - x_{g,c}\right)^2}\]

Skill Score -- относительное улучшение по отношению к
персистентному прогнозу (прогнозу «без изменений»):
\[\text{Skill} = 1 - \frac{\text{RMSE}_{model}}{\text{RMSE}_{persistence}}\]

ACC (Anomaly Correlation Coefficient) -- пространственная
корреляция аномалий:
\[\text{ACC} = \frac{\sum_{g,c} (x_{g,c} - \bar{x}_c)(\hat{x}_{g,c} - \bar{\hat{x}}_c)}{\sqrt{\sum_{g,c}(x_{g,c}-\bar{x}_c)^2 \cdot \sum_{g,c}(\hat{x}_{g,c}-\bar{\hat{x}}_c)^2}}\]

Все метрики вычислялись как для полного домена (61×41), так и для
центрального подрегиона Новосибирской области (53--57°с.ш., 79--85°в.д.,
425 узлов).

\subsection{Результаты экспериментов}\label{ux440ux435ux437ux443ux43bux44cux442ux430ux442ux44b-ux44dux43aux441ux43fux435ux440ux438ux43cux435ux43dux442ux43eux432}

Таблица 1 суммирует результаты основных конфигураций на горизонте +24 ч
для полного домена.

\textbf{Таблица 1.} Метрики качества на горизонте +24 ч, полный домен
61×41.

{\def\LTcaptype{none} % do not increment counter
\begin{longtable}[]{@{}lccc@{}}
\toprule\noalign{}
Конфигурация & RMSE↓ & Skill↑ & ACC↑ \\
\midrule\noalign{}
\endhead
\bottomrule\noalign{}
\endlastfoot
Персистентность (baseline) & 0.783 & 0\% & 0.377 \\
Базовая модель & 0.629 & 19.7\% & 0.451 \\
Только границы (bounds, \(b=5\)) & 0.589 & 32.2\% & 0.459 \\
Nudging, все переменные (\(\alpha=0.5\)) & 0.602 & 23.1\% & 0.540 \\
Nudging + границы & 0.562 & 35.2\% & 0.494 \\
ОИ, все перем. (\(L=300\) км, \(\sigma_o=0.5\)) & 0.538 & 38.0\% &
0.493 \\
ОИ + границы & \textbf{0.414} & \textbf{47.2\%} &
\textbf{0.649} \\
\end{longtable}
}

Таблица 2 -- результаты для регионального среза (Новосибирская область,
+24 ч).

\textbf{Таблица 2.} Метрики качества на горизонте +24 ч, регион НСКО
(53--57°с.ш., 79--85°в.д.).

{\def\LTcaptype{none} % do not increment counter
\begin{longtable}[]{@{}
  >{\raggedright\arraybackslash}p{(\linewidth - 6\tabcolsep) * \real{0.2105}}
  >{\centering\arraybackslash}p{(\linewidth - 6\tabcolsep) * \real{0.2632}}
  >{\centering\arraybackslash}p{(\linewidth - 6\tabcolsep) * \real{0.2632}}
  >{\centering\arraybackslash}p{(\linewidth - 6\tabcolsep) * \real{0.2632}}@{}}
\toprule\noalign{}
\begin{minipage}[b]{\linewidth}\raggedright
Конфигурация
\end{minipage} & \begin{minipage}[b]{\linewidth}\centering
RMSE↓
\end{minipage} & \begin{minipage}[b]{\linewidth}\centering
Skill↑
\end{minipage} & \begin{minipage}[b]{\linewidth}\centering
ACC↑
\end{minipage} \\
\midrule\noalign{}
\endhead
\bottomrule\noalign{}
\endlastfoot
Персистентность (baseline) & 0.660 & 0\% & 0.204 \\
Базовая модель & 0.621 & 15.5\% & 0.170 \\
Только границы & 0.621 & 15.5\% & 0.170 \\
Nudging, все переменные (\(\alpha=0.5\)) & 0.598 & 20.3\% & 0.198 \\
ОИ, все перем. (\(L=300\) км, \(\sigma_o=0.5\)) & \textbf{0.404} &
\textbf{40.1\%} & \textbf{0.306} \\
\end{longtable}
}

Анализ результатов. Наиболее значимые наблюдения:

\begin{enumerate}
\def\labelenumi{\arabic{enumi}.}
\item
  Сшивание границ само по себе даёт существенный прирост на
  полном домене (Skill с 19.7\% до 32.2\%) -- основной вклад идёт от
  исправления ошибок у краёв расчётной области. На центральном
  подрегионе (НСКО) этот эффект менее выражен -- Skill практически не
  меняется, что подтверждает локализацию граничных артефактов вблизи
  краёв домена.
\item
  Nudging даёт умеренное улучшение на обоих доменах. Его эффект
  сильнее при совместном использовании со сшиванием границ.
\item
  ОИ превосходит Nudging как на полном домене, так и особенно
  на центральном подрегионе (RMSE с 0.621 до 0.404, Skill с 15.5\% до 40.1\%).
  Это объясняется способностью ОИ распространять информацию от
  немногочисленных наблюдений на соседние узлы, тогда как Nudging
  напрямую корректирует только наблюдаемые узлы.
\item
  Комбинация ОИ + границы на полном домене даёт наилучший
  результат: Skill 47.2\%, ACC = 0.649 -- значительное превышение всех
  остальных конфигураций.
\end{enumerate}

\subsection{Анализ влияния состава усваиваемых переменных}\label{ux430ux43dux430ux43bux438ux437-ux432ux43bux438ux44fux43dux438ux44f-ux441ux43eux441ux442ux430ux432ux430-ux443ux441ux432ux430ux438ux432ux430ux435ux43cux44bux445-ux43fux435ux440ux435ux43cux435ux43dux43dux44bux445}

Таблица 3 показывает влияние выбора набора усваиваемых переменных (метод
Nudging, \(\alpha=0.5\), полный домен, +24 ч).

\textbf{Таблица 3.} Зависимость RMSE от набора усваиваемых переменных
(Nudging, +24 ч).

{\def\LTcaptype{none} % do not increment counter
\begin{longtable}[]{@{}llcc@{}}
\toprule\noalign{}
Набор & Переменные & RMSE↓ & Skill↑ \\
\midrule\noalign{}
\endhead
\bottomrule\noalign{}
\endlastfoot
Все 15 переменных & все & 0.602 & 23.1\% \\
Динамические & ветер, геопотенциал & 0.608 & 22.3\% \\
Приземные & t2m, msl, ветер & 0.614 & 21.6\% \\
Только температура & t2m, t850, t500 & 0.629 & 19.7\% \\
\end{longtable}
}

Усвоение только температуры практически не улучшает прогноз. Причина --
атмосферная динамика: если прогнозируемый ветер неверен,
скорректированная температура быстро «исправится» обратно динамическими
процессами за 1--2 шага. Для ОИ наблюдается та же тенденция: наилучшие
результаты при усвоении всех или динамических переменных.

\subsection{Подбор гиперпараметров}\label{ux43fux43eux434ux431ux43eux440-ux433ux438ux43fux435ux440ux43fux430ux440ux430ux43cux435ux442ux440ux43eux432}

Nudging -- выбор \(\alpha\). Эксперименты с
\(\alpha \in \{0.1, 0.5, 0.9, 1.3\}\) показали, что оптимальным является
\(\alpha = 0.5\). При \(\alpha > 0.9\) модель «перешагивает» и вносит
высокочастотные аномалии в поле, при \(\alpha < 0.3\) эффект коррекции
слишком мал.

ОИ -- выбор \(L\) и \(\sigma_o\). Результаты поиска по сетке
приведены в таблице 4.

\textbf{Таблица 4.} Подбор параметров ОИ: RMSE региона при горизонте +24
ч.

{\def\LTcaptype{none} % do not increment counter
\begin{longtable}[]{@{}cccc@{}}
\toprule\noalign{}
\(L\), км & \(\sigma_o = 0.2\) & \(\sigma_o = 0.5\) &
\(\sigma_o = 0.8\) \\
\midrule\noalign{}
\endhead
\bottomrule\noalign{}
\endlastfoot
50 & 0.534 & 0.499 & 0.497 \\
150 & 0.507 & 0.443 & 0.425 \\
300 & 0.414 & 0.404 & \textbf{0.402} \\
\end{longtable}
}

Лучшие результаты достигаются при \(L = 300\) км и
\(\sigma_o = 0.5\)--\(0.8\). Большой радиус корреляции обоснован тем,
что реальные пространственные корреляции ошибок нейросетевой модели
имеют крупный масштаб (несколько сотен километров), превышающий таковой
для традиционного NWP.

Ширина буфера \(b\). Тесты с \(b \in \{3, 5, 7\}\) показали,
что \(b = 5\) узлов (≈ 125 км) выгодно балансирует между шириной зоны
сопряжения и потерей «рабочей» площади прогноза.


\newpage

\section{МУЛЬТИРЕЗОЛЮЦИОННЫЙ ПОДХОД К ПРОГНОЗУ}\label{ux43cux443ux43bux44cux442ux438ux440ux435ux437ux43eux43bux44eux446ux438ux43eux43dux43dux44bux439-ux43fux43eux434ux445ux43eux434-ux43a-ux43fux440ux43eux433ux43dux43eux437ux443}

\subsection{Глобальная базовая модель высокого разрешения}\label{ux433ux43bux43eux431ux430ux43bux44cux43dux430ux44f-ux431ux430ux437ux43eux432ux430ux44f-ux43cux43eux434ux435ux43bux44c-ux432ux44bux441ux43eux43aux43eux433ux43e-ux440ux430ux437ux440ux435ux448ux435ux43dux438ux44f}

Для регионального дообучения используется глобальная модель версии v2,
обученная на ERA5 2010--2019 на сетке 512$\times$256 по 19 переменным.
Ключевые архитектурные улучшения по сравнению с прототипом v1 состояли в
следующем:
\begin{itemize}
\item замена простой графовой свёртки (GCN) в процессоре на
полноценный блок Interaction Network из 12 шагов message passing с
обновлением рёбер и узлов, что позволило модели учитывать взаимодействия
между узлами более гибко;
\item добавление 4-мерных признаков рёбер (расстояние и 3D-положение
концов в декартовых координатах), благодаря чему передача информации
между узлами стала явно опираться на сферическую геометрию Земли;
\item введены остаточные соединения на узлах и рёбрах, что ускорило
сходимость при увеличении глубины модели;
\item размерность скрытого состояния увеличена с 128 до 256;
\item активация во всех MLP заменена с PReLU на Swish (более
стабильная гладкая форма).
\end{itemize}
Итоговое число параметров -- 5.9~млн против 210~тыс у v1, при этом
время выдачи одного прогноза на GPU H100 остаётся порядка секунды.

\textbf{Таблица 5.} Результаты глобальной модели v2 (512×256, 19
переменных, 200 тестовых примеров).

{\def\LTcaptype{none} % do not increment counter
\begin{longtable}[]{@{}cccc@{}}
\toprule\noalign{}
Горизонт & RMSE (норм.) & Skill & ACC \\
\midrule\noalign{}
\endhead
\bottomrule\noalign{}
\endlastfoot
+6ч & 0.1577 & 57.1\% & 0.9875 \\
+12ч & 0.2074 & 59.7\% & 0.9782 \\
+18ч & 0.2533 & 57.7\% & 0.9669 \\
+24ч & 0.2901 & 55.5\% & 0.9563 \\
Среднее & 0.2321 & 57.3\% & 0.9722 \\
\end{longtable}
}

\subsection{Концепция мультирезолюционной сетки}\label{ux43aux43eux43dux446ux435ux43fux446ux438ux44f-ux43cux443ux43bux44cux442ux438ux440ux435ux437ux43eux43bux44eux446ux438ux43eux43dux43dux43eux439-ux441ux435ux442ux43aux438}

Классический подход к региональному нейросетевому прогнозированию
состоит в вырезании домена из глобальной модели, что неизбежно порождает
граничные артефакты и теряет информацию о соседних областях.
Предложенный мультирезолюционный подход объединяет глобальные и
региональные узлы в единую плоскую сетку без явных границ.

Сетка строится по следующему принципу:

\begin{enumerate}
\def\labelenumi{\arabic{enumi}.}
\tightlist
\item
  Берётся полная глобальная сетка 512×256 (131 072 узла).
\item
  Внутри ограничивающего прямоугольника ROI (регион Красноярска,
  54--58°с.ш., 91--95°в.д.) все глобальные узлы удаляются.
\item
  Вместо них вставляются региональные узлы на сетке 0.25° (61×41 = 2 501
  узел).
\item
  Итоговая сетка: 131 072 − \(N_{del}\) + 2 501 = 133 279
  узлов.
\end{enumerate}

Построение графа выполняется стандартным образом для объединённого
облака точек: пространственный индекс типа KD-дерева (бинарное дерево,
рекурсивно расщепляющее пространство по координатным осям и позволяющее
эффективно находить все точки в заданном радиусе) соединяет узлы без
каких-либо явных поправочных формул на границе. Неявное сглаживание
перехода обеспечивается 12 шагами message passing в процессоре.

Для обучения и применения модели в региональные узлы подставляется реальное
поле ERA5 высокого разрешения 0.25°, а в глобальные узлы --
то же поле ERA5, но с грубым шагом сетки 512$\times$256. Получается
сшитый датасет: внутри ROI -- мелкое разрешение,
снаружи -- крупное. Все дальнейшие эксперименты на
мультирезолюционной сетке проводятся именно на этом сшитом датасете.

\subsection{Сшивание глобальных и региональных данных}\label{ux441ux448ux438ux432ux430ux43dux438ux435-ux433ux43bux43eux431ux430ux43bux44cux43dux44bux445-ux438-ux440ux435ux433ux438ux43eux43dux430ux43bux44cux43dux44bux445-ux434ux430ux43dux43dux44bux445}

Процедура сшивания (далее -- merge) является жёсткой склейкой без
математических формул сглаживания на стыке. Формально:

\[x_{flat} = \text{concat}\!\left[\, x_{global}^{outside\,ROI}, \; x_{regional}^{inside\,ROI} \,\right]\]

Два множества узлов не пересекаются:
\(\text{supp}(x_{global}) \cap \text{supp}(x_{regional}) = \varnothing\)
по построению. Граничные эффекты (физический разрыв из-за несовпадения
разрешений) сглаживаются неявно за счёт message passing -- граничные
узлы соединены рёбрами и обмениваются признаками на каждом шаге
Processor.

Данный подход принципиально отличается от метода Davies {[}4{]}: метод Дэвиса 
применяется на этапе после получения прогноза, тогда
как мультирезолюционная склейка является частью входных данных для GNN.

\subsection{Стратегия дообучения с заморозкой процессора}\label{ux441ux442ux440ux430ux442ux435ux433ux438ux44f-ux434ux43eux43eux431ux443ux447ux435ux43dux438ux44f-ux441-ux437ux430ux43cux43eux440ux43eux437ux43aux43eux439-processor}

Проблема прямого дообучения глобальной модели на региональном датасете
-- катастрофическое забывание: процессор, содержащий основные
знания о динамике атмосферы, может деградировать при обучении на малом
региональном наборе.

Предложена двухфазная стратегия заморозки процессора: на первых
эпохах все параметры процессора фиксируются, а обучаются только
кодировщик и декодировщик, приспособливаясь к новой геометрии сетки
и статистике региона; затем процессор размораживается, но с
уменьшенным в 10~раз локальным темпом обучения, чтобы избежать
разрушения выученной глобальной динамики:

\begin{enumerate}
\def\labelenumi{\arabic{enumi}.}
\tightlist
\item
  Фаза~1 (6 эпох): процессор полностью заморожен, обучаются
  только кодировщик и декодировщик.
  \begin{itemize}
  \tightlist
  \item кодировщик адаптирует вложения к новой конфигурации сетки;
  \item декодировщик адаптирует проекцию к региональным статистикам.
  \end{itemize}
\item
  Фаза~2 (оставшиеся эпохи): процессор размораживается,
  локальный темп обучения для него уменьшается в 10~раз по сравнению
  с остальной частью модели.
\end{enumerate}

Сравнение с дообучением без заморозки (nofreeze) на 200 тестовых
примерах:

\textbf{Таблица 6.} Freeze6 vs nofreeze: метрики на тестовой выборке.

{\def\LTcaptype{none} % do not increment counter
\begin{longtable}[]{@{}lcc@{}}
\toprule\noalign{}
Метрика & freeze6 & nofreeze \\
\midrule\noalign{}
\endhead
\bottomrule\noalign{}
\endlastfoot
Skill (глобально, среднее +6\ldots+24ч) & \textbf{66.9\%} & 65.2\% \\
Skill (регион Красноярска, 45 узлов) & \textbf{75.8\%} & 74.5\% \\
ACC (глобально) & \textbf{0.983} & 0.981 \\
t2m RMSE в регионе, +24ч & \textbf{1.40°C} & 1.82°C \\
\end{longtable}
}

Более высокое значение t2m RMSE при nofreeze (+30\%) объясняется
деградацией процессора в зоне вне ROI, что ухудшает граничные условия и в
конечном счёте сказывается на точности в самом регионе.

\subsection{Результаты на регионе Красноярска}\label{ux440ux435ux437ux443ux43bux44cux442ux430ux442ux44b-ux43dux430-ux440ux435ux433ux438ux43eux43dux435-ux43aux440ux430ux441ux43dux43eux44fux440ux441ux43aux430}

\textbf{Таблица 7.} Мультирезолюционная модель (freeze6, 200
примеров): метрики в зоне Красноярска (45 внутренних узлов,
55.5--56.5°с.ш., 92--94°в.д.). RMSE приведён в нормированном виде
(после стандартизации по каналу), базовая RMSE рассчитывается для
персистентного прогноза.

{\def\LTcaptype{none} % do not increment counter
\begin{longtable}[]{@{}ccccc@{}}
\toprule\noalign{}
Горизонт & RMSE (норм.) & base\_RMSE & Skill & ACC \\
\midrule\noalign{}
\endhead
\bottomrule\noalign{}
\endlastfoot
+6ч & 0.0810 & 0.2243 & 63.9\% & 0.748 \\
+12ч & 0.1084 & 0.3919 & 72.3\% & 0.685 \\
+18ч & 0.1351 & 0.5322 & 74.6\% & 0.643 \\
+24ч & 0.1603 & 0.6351 & 74.8\% & 0.603 \\
Среднее & 0.1212 & 0.4459 & 73.9\% & 0.670 \\
\end{longtable}
}

Физические RMSE по ключевым переменным (горизонт +24ч):

{\def\LTcaptype{none} % do not increment counter
\begin{longtable}[]{@{}lclc@{}}
\toprule\noalign{}
Переменная & RMSE & Переменная & RMSE \\
\midrule\noalign{}
\endhead
\bottomrule\noalign{}
\endlastfoot
t2m & 1.40°C & u@850 & 2.55 м/с \\
msl & 2.34 Па & v@850 & 1.77 м/с \\
10u & 0.87 м/с & z@500 & 1.7 гпм \\
\end{longtable}
}

Для сравнения: автономная региональная GNN (раздел 6, регион НСКО)
давала t2m RMSE = 3.52°C при горизонте +24ч. Мультирезолюционный подход
с фризом процессора даёт 1.40°C -- улучшение в 2.5 раза за счёт
использования глобального контекста.

Замечание о различных геометриях оценки. В таблице~7
используется \emph{внутренняя зона} из 45~узлов (55.5--56.5°с.ш.,
92--94°в.д.), фокусирующаяся на ближайших окрестностях
Красноярска и пригодная для сравнения с пилотной региональной GNN
(раздел~6). В подразделе~7.7 и в апробационных
таблицах раздела~10 рассматривается \emph{расширенный ROI} из
2~501 узла (50--60°с.ш., 83--98°в.д.), совпадающий с ограничивающим
прямоугольником мультирезолюционной сетки и используемый для всех
экспериментов по DA. Базовые числа для двух геометрий ожидаемо
различаются (расширенный ROI включает пограничные узлы, где
ошибка выше), поэтому прямое сравнение абсолютных значений между
разделами некорректно -- следует сопоставлять приросты
относительно baseline.

\subsection{Дообучение на композитном датасете с остаточной связью}

Production-модель дообучается непосредственно на сшитом
датасете, где в узлах ROI используется реальное
поле ERA5 0.25°, а вне ROI -- ERA5 с грубой сеткой 512$\times$256.
Ключевые особенности конфигурации дообучения, добавленные по
сравнению с базовой стратегией заморозки процессора (раздел~7.4):

\begin{itemize}
\tightlist
\item
  модель предсказывает приращение поля относительно
  предыдущего шага, а не абсолютные
  значения, что снижает амплитуду ошибки на ближнем горизонте;
\item
  длина фазы заморозки процессора увеличена с 6 до 10 эпох --
  кодировщик и декодировщик дольше адаптируются к статистике
  реальных региональных полей;
\item
  снижен локальный темп обучения процессора при его разморозке;
\item
  базовый темп обучения уменьшен до $5\cdot10^{-5}$;
\item
  применяется учебный план (\emph{curriculum}) по длине прогноза:
  в первых эпохах модель предсказывает один шаг, затем длина
  rollout постепенно увеличивается до трёх шагов; это типичный
  приём борьбы с накоплением ошибки в авторегрессионных моделях.
\end{itemize}

Эта модель используется как базовая во всех экспериментах по
усвоению данных на мультирезолюционной сетке (раздел~7.7).

\subsection{Усвоение данных на мультирезолюционной модели}\label{ux443ux441ux432ux43eux435ux43dux438ux435-ux434ux430ux43dux43dux44bux445-ux43dux430-ux43cux443ux43bux44cux442ux438ux440ux435ux437ux43eux43bux44eux446ux438ux43eux43dux43dux43eux439-ux43cux43eux434ux435ux43bux438}

Для мультирезолюционной модели freeze6 проведён полный
вычислительный эксперимент по усвоению данных на сшитом датасете.
Прогонялась полная сетка из 8~значений радиуса корреляции
\(\times\) 3 значений \(\sigma_o\) \(\times\) 2~плотности наблюдений
(10\% и 1\%); дополнительно проведены эксперименты с Nudging при
5~значениях \(\alpha\) и сравнение для 4~групп усваиваемых переменных.
Параллельно проведена независимая серия DA-экспериментов на исходной
глобальной сетке 512$\times$256 (см.~подраздел~7.7.5).

\subsubsection{Подбор гиперпараметров ОИ}

\textbf{Таблица 8a.} ОИ при 10\% наблюдений, все 19 переменных: Skill
Score по региону Красноярска (+6~ч) на сшитом датасете (baseline
+6~ч = 59.42\%). Прочерки соответствуют комбинациям, для которых не
проводился полный sweep по $\sigma_o$: радиусы 25, 150, 200, 300,
500~км исследовались только при $\sigma_o = 0.5$, поскольку из
sweep'а на 10/50/100~км уже было видно, что $\sigma_o = 0.5$ --
оптимум, а изменение радиуса корреляции даёт более крупный эффект,
чем изменение $\sigma_o$.

{\def\LTcaptype{none}
\begin{longtable}[]{@{}cccc@{}}
\toprule\noalign{}
Радиус корр. $L$, км & $\sigma_o = 0.3$ & $\sigma_o = 0.5$ & $\sigma_o = 1.0$ \\
\midrule\noalign{}
\endhead
\bottomrule\noalign{}
\endlastfoot
10 & -- & 61.73\% & 60.93\% \\
25 & -- & 67.16\% & -- \\
50 & 76.39\% & 75.40\% & 71.32\% \\
100 & 77.77\% & \textbf{77.86\%} & 75.86\% \\
150 & -- & 76.83\% & -- \\
200 & -- & 75.62\% & -- \\
300 & -- & 73.84\% & -- \\
500 & -- & 71.62\% & -- \\
\end{longtable}
}

Оптимальная конфигурация: \(L = 100\) км, \(\sigma_o = 0.5\) -- Skill +6~ч
77.86\% (+18.4 п.п. к baseline). При усреднении по всем горизонтам
прогноза средний Skill составляет 82.83\% против 66.88\% у baseline
(см. таблицу~8c). Радиус 100~км соответствует масштабу мезо-\(\alpha\)
синоптических структур.

\subsubsection{Влияние состава усваиваемых переменных}

\textbf{Таблица 8b.} Зависимость Skill от набора усваиваемых переменных
(ОИ, 10\%, \(L = 100\) км, \(\sigma_o = 0.5\), +6~ч).

{\def\LTcaptype{none}
\begin{longtable}[]{@{}lcc@{}}
\toprule\noalign{}
Группа переменных & Skill & $\Delta$ к baseline \\
\midrule\noalign{}
\endhead
\bottomrule\noalign{}
\endlastfoot
Только t2m & 63.44\% & +4.0~п.п. \\
Приземные (7 пер.) & 64.49\% & +5.1~п.п. \\
Приземные + верхние (17 пер.) & 70.23\% & +10.8~п.п. \\
Все 19 переменных & \textbf{77.86\%} & \textbf{+18.4~п.п.} \\
\end{longtable}
}

Результат подтверждает выводы пилотных экспериментов:
усвоение только температуры практически бесполезно. Включение всех
переменных критически важно; при этом верхние уровни (850/500~гПа)
вносят основной вклад (+7– 10~п.п. сверх приземных).

\subsubsection{Сравнение Nudging и ОИ}

При 10\% наблюдений лучший Nudging (\(\alpha = 0.9\)) даёт
Skill 61.3\%, что значительно уступает ОИ (77.86\%).

При 1\% наблюдений (25 узлов) полный sweep по радиусу корреляции
представлен в таблице~8b' ниже. Для малого числа наблюдений
оптимальный радиус увеличивается (150--200~км вместо 100~км),
что физически обосновано необходимостью <<дотягивания>> информации
на большие расстояния.

\textbf{Таблица 8b'.} ОИ при 1\% наблюдений (≈25 узлов),
$\sigma_o = 0.5$, все 19 переменных: Skill +6~ч на сшитом датасете.

{\def\LTcaptype{none}
\begin{longtable}[]{@{}ccc@{}}
\toprule\noalign{}
Радиус корр. $L$, км & Skill +6~ч & ACC \\
\midrule\noalign{}
\endhead
\bottomrule\noalign{}
\endlastfoot
10  & 59.66\% & 0.9608 \\
50  & 62.08\% & 0.9625 \\
100 & 65.67\% & 0.9630 \\
150 & 67.52\% & 0.9619 \\
200 & \textbf{67.88\%} & 0.9604 \\
300 & 67.38\% & 0.9577 \\
500 & 66.86\% & 0.9569 \\
\end{longtable}
}

Оптимум при 1\% наблюдений: $L = 200$~км, Skill +6~ч 67.88\%
(+8.5~п.п. к baseline 59.42\%). По горизонтам прогноза:
+12~ч 74.25\%, +18~ч 76.59\%, +24~ч 77.86\%; t2m RMSE в регионе
растёт от 1.22°C (+6~ч) до 1.79°C (+24~ч). При сравнении
с 10\% наблюдений видно, что десятикратное уменьшение плотности
сети снижает выигрыш от ОИ примерно вдвое (с +18.4 до +8.5~п.п.
на +6~ч), но эффект остаётся значимым.

\subsubsection{Динамика по горизонтам прогноза}

\textbf{Таблица 8c.} Skill Score по горизонтам прогноза при оптимальной
ОИ (все 19 пер., \(L = 100\) км, \(\sigma_o = 0.5\), 10\%).

{\def\LTcaptype{none}
\begin{longtable}[]{@{}ccc@{}}
\toprule\noalign{}
Горизонт & Baseline & ОИ \\
\midrule\noalign{}
\endhead
\bottomrule\noalign{}
\endlastfoot
+6~ч & 59.42\% & \textbf{77.86\%} \\
+12~ч & 68.95\% & \textbf{82.00\%} \\
+18~ч & 73.48\% & \textbf{83.54\%} \\
+24~ч & 76.21\% & \textbf{84.28\%} \\
Среднее & 66.88\% & \textbf{82.83\%} \\
\end{longtable}
}

Максимальный абсолютный прирост наблюдается на ближних горизонтах
(+18.4~п.п. при +6~ч), где усвоение непосредственно корректирует
текущее состояние. С увеличением горизонта прогноза относительный
эффект DA уменьшается, но остаётся значительным (+8.1~п.п. при +24~ч).

Основные выводы по DA на мультирезолюционной модели:
\begin{enumerate}
\def\labelenumi{\arabic{enumi}.}
\tightlist
\item ОИ стабильно превосходит Nudging при любой плотности наблюдений.
\item Оптимальный радиус корреляции \(L = 100\) км; при малом числе
  наблюдений (1\%) оптимум сдвигается к 150--200~км.
\item Дообучение на сшитом датасете даёт модели наибольший резерв для
  улучшения путём усвоения (+15.95~п.п. в среднем по горизонтам).
\item Необходимо усваивать все переменные, включая верхние уровни
  атмосферы.
\end{enumerate}

\subsubsection{Усвоение данных на глобальной сетке 512×256}\label{ux443ux441ux432ux43eux435ux43dux438ux435-ux43dux430-ux433ux43bux43eux431ux430ux43bux44cux43dux43eux439-ux441ux435ux442ux43aux435}

Параллельно с экспериментами на мультирезолюционной сетке проведена
независимая серия DA-экспериментов на исходной глобальной сетке
512$\times$256, используемой базовой моделью v2 (см.\,подраздел~7.1). Цель -- проверить, даёт ли
ОИ значимый прирост, если применять её напрямую к выходу глобальной
модели в зоне ROI Красноярска, без перехода на мультирезолюционную
сетку. Наблюдения вводятся только в узлы, попавшие в ROI
(294 узла на грубой сетке), 200 тестовых примеров; для оценки
качества используется та же геометрия ROI.

\textbf{Таблица 8d.} ОИ при 10\% наблюдений (29 станций),
$\sigma_o = 0.5$: Skill региона на глобальной сетке (baseline +6~ч
77.77\%, +24~ч 78.20\%).

{\def\LTcaptype{none}
\begin{longtable}[]{@{}cccc@{}}
\toprule\noalign{}
$L$, км & +6~ч Skill & +24~ч Skill & ACC (+6~ч) \\
\midrule\noalign{}
\endhead
\bottomrule\noalign{}
\endlastfoot
50  & 79.36\% & 80.37\% & 0.8930 \\
100 & 80.78\% & 83.48\% & 0.8976 \\
150 & \textbf{81.17\%} & 85.21\% & \textbf{0.9006} \\
200 & 80.99\% & \textbf{85.81\%} & 0.9018 \\
300 & 80.43\% & 85.62\% & 0.9016 \\
500 & 79.96\% & 84.80\% & 0.8999 \\
\end{longtable}
}

Оптимальный радиус корреляции на глобальной сетке сдвинут к
150--200~км (вместо 100~км на мультирезолюционной сетке 0.25°),
что объясняется большим шагом сетки -- эффективная связность узлов
требует большего радиуса. Прирост Skill при +24~ч составляет
+7.6~п.п. (с 78.20\% до 85.81\%); при 1\% наблюдений (3 станции)
выигрыш падает до $\le 0.3$~п.п., что подтверждает критическую
зависимость качества DA от плотности сети наблюдений. Абсолютные
значения Skill на глобальной сетке выше, чем на мультирезолюционной,
поскольку метрика рассчитывается на грубой сетке 0.703°, где
потенциальная мелкомасштабная изменчивость отсутствует.


\newpage

\section{КАСКАДНЫЙ НЕЙРОСЕТЕВОЙ DOWNSCALER}\label{ux43aux430ux441ux43aux430ux434ux43dux44bux439-ux43dux435ux439ux440ux43eux441ux435ux442ux435ux432ux43eux439-downscaler}

\subsection{Концепция каскадного downscaling}\label{ux43aux43eux43dux446ux435ux43fux446ux438ux44f-ux43aux430ux441ux43aux430ux434ux43dux43eux433ux43e-downscaling}

GNN-модель оперирует на грубой сетке 0.25° и неизбежно «размазывает»
локальные особенности рельефа и городского микроклимата. Для их
восстановления предлагается каскадная схема: GNN даёт поле
первого приближения, которое затем уточняется специализированной сетью
WeatherUNet, обучаемой предсказывать поправку к полю первого
приближения:

\[\hat{x}_{fine} = \hat{x}_{GNN} + \Delta(z_{surf},\, \hat{x}_{GNN}), \quad \Delta = \text{UNet}(\hat{x}_{GNN},\, z_{surf},\, \text{lsm})\]

Концептуально это соответствует традиционному статистическому
даунскейлингу {[}16{]}, но реализованному нейросетью.

\subsection{Архитектура WeatherUNet}\label{ux430ux440ux445ux438ux442ux435ux43aux442ux443ux440ux430-weatherunet}

Использована архитектура типа кодировщик--декодировщик с
симметричными перекрёстными связями (U-Net {[}17{]}). В качестве
итоговой конфигурации после экспериментов выбрана модель WeatherUNet
с расширениями, подавляющими сглаживание поля и улучшающими
точность на ближнем горизонте:

\begin{itemize}
\tightlist
\item
  3-уровневый кодировщик из блоков ResConvBlock c
  Squeeze--Excitation;
\item
  блок self-attention на узком месте кодировщика для учёта дальних
  пространственных связей;
\item
  спектральные свёртки в декодировщике для
  восстановления высокочастотных компонент поля;
\item
  окно предыстории 24~ч (4~временных среза). Входные каналы:
  $4 \times 19$ динамических + 2 статических, выходные -- 19~полей;
\item
  функция потерь -- MSE с дополнительными спектральной и градиентной
  составляющими (веса 0.1 и 0.05) для подавления излишнего
  сглаживания;
\item
  оптимизатор Adam с расписанием OneCycleLR, 80 эпох обучения;
\item
  число параметров: ~25\,500\,000.
\end{itemize}

Сеть предсказывает не абсолютное значение тонкого поля, а
поправку к полю первого приближения
$\Delta = \hat{x}_{fine} - \hat{x}_{coarse}$. Такая постановка
упрощает задачу обучения: вместо восстановления полного поля сети
достаточно выучить локальные поправки, обусловленные рельефом и
статистикой региональных полей.

\subsection{Роль WeatherUNet в каскаде}\label{ux430ux432ux442ux43eux43dux43eux43cux43dux44bux439-unet}

WeatherUNet рассматривается исключительно как уточняющий блок в каскаде
с мультирезолюционной GNN, а не как самостоятельная прогностическая
модель. Локальная свёрточная архитектура без доступа к синоптической
динамике не способна устойчиво развёртывать прогноз на 24~ч вперёд;
её сильная сторона~-- восстановление мелкомасштабных пространственных
структур (орографические эффекты, мезомасштабные особенности осадков и
давления) поверх уже физически согласованного поля GNN. Поэтому на
каждом шаге авторегрессии на вход WeatherUNet подаётся свежий прогноз
GNN, а сеть выучивает только поправку $\Delta$ к этому полю. Полные
результаты каскада приведены ниже в разделе~8.4.

\subsection{Результаты полного каскада GNN и UNet}\label{ux440ux435ux437ux443ux43bux44cux442ux430ux442ux44b-ux43fux43eux43bux43dux43eux433ux43e-ux43aux430ux441ux43aux430ux434ux430-gnnunet}

Полный каскад применяет WeatherUNet к прогнозу мультирезолюционной
GNN на каждом шаге авторегрессии:
$\hat{x}_{fine} = \hat{x}_{GNN} + \text{UNet}(\hat{x}_{GNN}, z_{surf}, \text{lsm})$.
Оценка проведена на ERA5 0.25° (200 примеров, AR=4) относительно
базовой GNN-модели без каскада.

\textbf{Таблица 9.} Каскад против чистого GNN: RMSE по температуре
у поверхности (200 примеров, ERA5 0.25°, физические единицы).

{\def\LTcaptype{none} % do not increment counter
\begin{longtable}[]{@{}lccc@{}}
\toprule\noalign{}
Горизонт & GNN, °C & Каскад, °C & $\Delta$, \% \\
\midrule\noalign{}
\endhead
\bottomrule\noalign{}
\endlastfoot
+6~ч  & 1.67 & 1.69 & $-1.2$ \\
+12~ч & 1.95 & 1.97 & $-1.4$ \\
+18~ч & 2.11 & 2.14 & $-1.2$ \\
+24~ч & 2.24 & 2.27 & $-1.4$ \\
\end{longtable}
}

Ключевые наблюдения:

\begin{itemize}
\tightlist
\item
  каскад не уступает чистому GNN по температуре~-- разница в
  RMSE не превышает 1.4~\% во всём диапазоне +6...+24~ч и
  лежит существенно ниже типичной ошибки прогноза;
\item
  по приземному давлению (msl) каскад улучшает прогноз на
  1--5~\%, по геопотенциалу на 850~гПа~-- на 2--3~\%;
\item
  по осадкам (tp) каскад даёт качественный прирост: skill
  относительно базовой GNN растёт на 80--85~\% на всех горизонтах~--
  UNet восстанавливает мезомасштабные структуры осадков, сглаженные
  глобальной сеткой;
\item
  в среднем по всем 19 переменным изменение RMSE укладывается в
  $\pm 1$~\%~-- каскад нейтрален к интегральному качеству и при
  этом заметно улучшает прогноз ключевых для прикладных задач полей
  (осадки, давление).
\end{itemize}

Вывод. Каскад GNN и WeatherUNet не уступает мультирезолюционной
GNN по температуре и значительно улучшает прогноз осадков и
давления. На гладких термодинамических полях, где мезомасштабная
структура слаба (температура и геопотенциал на 500~гПа), каскад
скорее не нужен, но и не вредит. В оперативной системе UNet
включён в пайплайн на всех каналах, поскольку выигрыш на осадках
существенно превышает пренебрежимо малую потерю на температуре.


\newpage

\section{ПОСТПРОЦЕССИНГ ПРОГНОЗА}\label{ux43fux43eux441ux442ux43fux440ux43eux446ux435ux441ux441ux438ux43dux433-ux43fux440ux43eux433ux43dux43eux437ux430}

\subsection{Проблема репрезентативности ERA5 в городах}\label{ux43fux440ux43eux431ux43bux435ux43cux430-ux440ux435ux43fux440ux435ux437ux435ux43dux442ux430ux442ux438ux432ux43dux43eux441ux442ux438-era5-ux432-ux433ux43eux440ux43eux434ux430ux445}

Тестирование на ERA5 даёт хорошие метрики, потому что и обучение,
и валидация ведутся на одних и тех же сглаженных полях реанализа.
При переходе к верификации по данным городских метеостанций качество
систематически падает~-- это явление известно как ошибка
репрезентативности {[}18{]}.

Основные причины несоответствия поля ERA5 и станционных наблюдений:

\begin{enumerate}
\def\labelenumi{\arabic{enumi}.}
\tightlist
\item
  Усреднение по ячейке сетки. Ячейка ERA5 (\textasciitilde31~км)
  усредняет внутри себя городскую застройку, реку, промышленную зону и
  пригород. Тонкие вертикальные структуры и горизонтальные
  неоднородности, определяющие погоду в конкретной точке, теряются
  {[}19{]}.
\item
  Городской тепловой остров. Тепловые острова города обусловлены
  теплоёмкостью материалов покрытий, антропогенным тепловыделением и
  стеснённостью уличных каньонов. ERA5 не разрешает эти эффекты,
  что особенно заметно в ночные часы {[}20{]}.
\item
  Сезонная зависимость смещения. Смещение ERA5 относительно WMO~29570
  (Красноярск, ст. Центральная) существенно меняется в зависимости от
  сезона и времени суток: зимой станция теплее ERA5 на 5--7~°C (вклад
  городского острова и температурных инверсий), летом днём ERA5 теплее
  на 3--6~°C (перегрев сглаженной ячейки).
\end{enumerate}

Красноярск представляет особый случай: Енисей не замерзает
зимой (тепловой контраст), частые инверсии и сложный рельеф
дополнительно усиливают эффект. Спутниковый тепловой остров
подтверждён по данным Landsat {[}21{]}.

\subsection{Поправка на разницу высот по вертикальному градиенту температуры}\label{ux43fux43eux43fux440ux430ux432ux43aux430-ux432ux435ux440ux442ux438ux43aux430ux43bux44cux43dux44bux439-ux433ux440ux430ux434ux438ux435ux43dux442}

Высота модельного узла сетки (усреднённый рельеф ячейки ERA5) и высота
реальной станции могут существенно отличаться. Поправка вычисляется по
стандартному вертикальному градиенту температуры тропосферы:

\[T_{corr} = T_{model} + \Gamma \cdot (z_{station} - z_{grid}), \quad \Gamma = 6.5 \text{ °C/км}\]

где \(z_{station}\) -- высота целевой точки над уровнем моря,
\(z_{grid}\) -- высота ближайшего узла сетки по полю \(z_{surf}\) ERA5
(\(z_{surf}/g\), где \(g = 9.80665\) м/с²).

Для 19 станций в окрестности Красноярска разница высот составляет от −30
до +80 м, что даёт поправку порядка ±0.2--0.5°C. Метод не имеет
свободных параметров и применяется первым в пайплайне.

\subsection{Learned MOS -- статистическая постобработка
модельных результатов}\label{learned-mos-ux441ux442ux430ux442ux438ux441ux442ux438ux447ux435ux441ux43aux430ux44f-ux43fux43eux441ux442ux43eux431ux440ux430ux431ux43eux442ux43aux430-ux43cux43eux434ux435ux43bux44cux43dux44bux445-ux440ux435ux437ux443ux43bux44cux442ux430ux442ux43eux432}

Model Output Statistics (MOS) {[}9{]} -- классический метод
статистической постобработки: регрессионная модель обучается
предсказывать систематическую расстройку (bias) между прогнозом и
наблюдением с учётом конкретного синоптического контекста.

Для каждой из 19 станций (NOAA ISD-Lite, 2016--2024) обучается отдельная
регрессия на базе гистограммного градиентного бустинга
(\texttt{HistGradientBoostingRegressor}), предсказывающая:

\[\widehat{\text{bias}}(g, c, t) = T_{station} - T_{model}, \quad \text{bias} = f(\mathbf{x}_{20})\]

Вектор признаков состоит из 20 компонент, объединённых в пять смысловых групп.
\textbf{Метеорологическая обстановка} (8 признаков) описывает текущее
состояние атмосферы в узле сетки модели: температура воздуха на 2~м,
температура точки росы, модуль приземного ветра и его направление,
представленное синусом и косинусом (чтобы избежать разрыва на 360°),
приземное давление, интегральная влажность атмосферного столба и
суммарные осадки за 6 часов.
\textbf{Временной контекст} (5 признаков): час суток и день года,
каждый закодирован парой «синус/косинус» (циклическое представление,
избавляющее модель от искусственного разрыва между 23:00 и 00:00),
плюс высота солнца над горизонтом. Эта группа отвечает за сезонные и
суточные систематические ошибки модельного поля.
\textbf{Лаговые признаки} (2): значение температуры в том же узле
6 часов назад и её приращение за этот же интервал. Они передают модели MOS
краткосрочную динамику в районе станции.
\textbf{Пространственное положение станции} (3): её широта, долгота и
высота над уровнем моря.
\textbf{Радиационная обстановка} (2): мощность приходящей коротковолновой
радиации и общая облачность. Эти два признака приближённо описывают,
насколько активно поверхность прогревается солнцем и охлаждается
за счёт длинноволнового излучения ночью~-- то есть содержат ключевую
физическую информацию для коррекции отклонения ERA5 от реальной точки.

Метрика на тестовом периоде (2024 год): MAE = 1.32~°C, что
существенно лучше выхода ERA5 без постпроцессинга, где средняя
абсолютная ошибка относительно станций составляет 3--5~°C в зависимости
от сезона.

\subsection{Spatial IDW -- пространственное распределение поправки}\label{spatial-idw-ux43fux440ux43eux441ux442ux440ux430ux43dux441ux442ux432ux435ux43dux43dux43eux435-ux440ux430ux441ux43fux440ux435ux434ux435ux43bux435ux43dux438ux435-ux43fux43eux43fux440ux430ux432ux43aux438}

Learned MOS корректирует прогноз лишь в 19 точках. Для распространения
поправки на все 2~501 узлов региональной сетки используется метод
обратного взвешения расстояния (IDW, Шепард {[}22{]}):

\[\text{bias}(v) = \frac{\displaystyle\sum_{k:\, d_{vk} < r_{max}} w_k \cdot \text{bias}(s_k)}{\displaystyle\sum_{k:\, d_{vk} < r_{max}} w_k}, \quad w_k = d(v, s_k)^{-p}\]

Используются следующие параметры: степень \(p = 2\) (обратный
квадрат расстояния), максимальный радиус поиска \(r_{max} = 300\)~км;
расстояния вычисляются по формуле гаверсинуса.

Если в радиусе \(r_{max}\) нет ни одной станции, поправка для данного
узла принимается равной нулю. Метод гарантирует \(C^0\)-непрерывность
поля поправок и сохраняет точные значения MOS в узлах, совпадающих со
станциями.

Итоговый порядок применения постпроцессинга следующий:
выход GNN последовательно проходит (1)~поправку на разницу высот
по вертикальному градиенту, (2)~Learned MOS в станционных точках
и (3)~распространение полученной поправки на все узлы регионального
графа методом IDW. В оперативной схеме дополнительно применяется
коррекция систематического смещения самих полей GDAS перед подачей
в модель.

\paragraph{Границы применимости.}
Описанная последовательность поправок осмысленна только тогда, когда
конечный потребитель прогноза~-- набор реальных метеостанций.
Если же качество оценивается на сетке ERA5 или GDAS, цепочка
коррекций ухудшает RMSE: станционные значения и реанализ
описывают разные физические объекты~-- точечное наблюдение в городе
против усреднения по ячейке размером порядка 30~км. В контрольной
верификации на ERA5 в окрестности Красноярска RMSE по $t2m$ растёт с
2.08~°C (исходный выход GNN) до 6.49~°C после полного применения
последовательности lapse-rate, MOS и IDW. В оперативном режиме эта
проблема не возникает автоматически: прогноз отдаётся в координатах
конкретных станций.

\subsection{Сезонное масштабирование скорости ветра}\label{ux441ux435ux437ux43eux43dux43dux43eux435-ux43cux430ux441ux448ux442ux430ux431ux438ux440ux43eux432ux430ux43dux438ux435-ux432ux435ux442ux440ux430}

ERA5 систематически завышает приземную скорость ветра относительно
наблюдений ISD-Lite (NOAA Integrated Surface Database -- открытая база
часовых метеоданных мировой сети станций ВМО, поддерживаемая Национальным
центром климатических данных NOAA NCEI~[23]). По выборке из 19~станций
Красноярского края за 2022--2024~гг. среднее завышение составляет
$\sim$27\%. При этом GDAS/GFS\footnote{Global Data Assimilation System и
Global Forecast System -- оперативный анализ и прогностическая модель,
регулярно публикуемые NOAA NCEP с разрешением 0.25° каждые 6~часов; в
отличие от ERA5, доступны в режиме, близком к реальному времени.~[24]},
которые используются для оперативной инициализации, имеют собственное
смещение относительно ERA5, причём это смещение меняет знак по сезонам
(+15\% весной, $-6$\% зимой). Совместное действие двух систематических
ошибок делает простой статический MOS для скорости ветра неэффективным.
Для температуры аналогичной проблемы не возникает, поскольку модель
Learned MOS (раздел~9.3) в число входных признаков включает синус и
косинус дня года, час суток и lag6h~-- HistGBR-регрессор фактически
обучает свой собственный сезонно-зависимый отклик и сам ловит
разнознаковые систематические ошибки. Для ветра такой ML-модели
не обучено; вместо этого предложена прозрачная мультипликативная
коррекция с помесячным разрешением:

\[
\text{wind}_{\text{corr}}(s, m) = \text{wind}_{\text{model}}(s, m) \cdot
\frac{\overline{\text{ISD}}_m(s)}{\overline{\text{ERA5}}_m(s)},
\]

где $s$~-- индекс станции, $m \in \{1, \dots, 12\}$~-- календарный
месяц, а средние $\overline{\text{ISD}}_m$ и $\overline{\text{ERA5}}_m$
вычисляются один раз по совместной выборке 2022--2024~гг. с раздельной
фильтрацией по дневным и ночным часам. Полученные таблицы коэффициентов
сохраняются и применяются на этапе вывода прогноза: никакого повторного
обучения модели по месяцам не происходит~-- каждое значение коэффициента~--
это просто заранее посчитанная статистика отношения наблюдений к
реанализу для конкретной станции и конкретного месяца. Сами по себе
эти коэффициенты несут физический смысл локального масштабного фактора:
например, значение $0.34$ для январских ветров в Емельяново означает,
что реальный средний январский ветер на этой станции составляет около
34\% от того, что выдаёт ERA5,~-- очень крупное смещение, типичное для
долины Енисея в условиях зимних инверсий и слабой циркуляции.

\textbf{Таблица 10.} Месячные коэффициенты $\text{ISD}/\text{ERA5}$
для скорости ветра (выдержка по основным станциям Красноярского края).

{\def\LTcaptype{none}
\begin{longtable}[]{@{}lccccc@{}}
\toprule\noalign{}
Станция & Январь & Апрель & Июль & Октябрь & Год \\
\midrule\noalign{}
\endhead
\bottomrule\noalign{}
\endlastfoot
Емельяново & 0.34 & 1.24 & 0.98 & 0.72 & 0.90 \\
Ачинск & 0.63 & 0.65 & 0.81 & 0.66 & 0.66 \\
Канск & 0.35 & 0.68 & 0.54 & 0.52 & 0.53 \\
\end{longtable}
}

Реализация. Коэффициенты, полученные в станционных
координатах, распространяются на все узлы регионального грида
посредством IDW-интерполяции (\(p=2\), \(r_{\max}=300\)~км). Для
устойчивости вводятся два защитных порога: минимальная скорость ветра
$0.3$~м/с (снизу), при которой коррекция не применяется, и нижний
предел самого коэффициента $0.2$ (защита от деления на близкие к нулю
ERA5-средние). Поправка применяется к зональной ($u_{10}$, направленной
с запада на восток) и меридиональной ($v_{10}$, направленной с юга на
север) компонентам приземного ветра.

\subsection{Вычисление скорости и направления
ветра}\label{wind-speed-direction}

Модель предсказывает не саму скорость и направление приземного ветра,
а две его декартовы компоненты на высоте 10~метров над поверхностью:
зональную $u_{10}$ (положительное значение~-- ветер с запада на восток)
и меридиональную $v_{10}$ (положительное значение~-- ветер с юга на север);
нижний индекс ``10'' соответствует стандартной для синоптики высоте
приземного ветра. Производные характеристики вычисляются на этапе
парсинга прогноза.

Скорость ветра ($\text{speed}$) и его направление
($\text{dir}$, метеорологический угол, отсчитываемый от направления, откуда
дует ветер):
\[
  \text{speed} = \sqrt{u_{10}^2 + v_{10}^2}, \qquad
  \text{dir} = \bigl(270° - \arctan2(v_{10},\, u_{10})\bigr) \bmod 360°.
\]

Оценка порывов ветра $\text{gust}$ строится по эмпирической формуле,
связывающей приземную скорость со скоростью на уровне 850~гПа
(нижняя тропосфера) и наглядно отражающей вклад вертикального
перемешивания:
\[
  \text{gust} = \text{speed}_{10\text{m}} + 0{,}6 \cdot
  \max\!\bigl(\text{speed}_{850} - \text{speed}_{10\text{m}},\; 0\bigr).
\]

Пространственное усреднение по зонам.
Для формирования сводки по зонам (центр города -- 3 точки, город +
окрестности -- 9 точек, регион -- 45 точек) используется
\emph{векторное} среднее: усредняются компоненты $u$ и $v$, после чего
вычисляется скорость из средних компонент. Это физически корректно,
поскольку арифметическое среднее скоростей завышает результат при
различающихся направлениях.

Тип осадков определяется по температуре: $t \leq 1$~°C -- снег,
$1$~°C $< t \leq 3$~°C -- мокрый снег, $t > 3$~°C -- дождь. Порог
значимости осадков: 0.1~мм/6ч.


\newpage

\section{АПРОБАЦИЯ В ОПЕРАТИВНОМ РЕЖИМЕ}\label{ux430ux43fux440ux43eux431ux430ux446ux438ux44f-ux432-ux43eux43fux435ux440ux430ux442ux438ux432ux43dux43eux43c-ux440ux435ux436ux438ux43cux435}

\subsection{Конфигурация оперативного прогноза по данным GDAS}\label{ux43aux43eux43dux444ux438ux433ux443ux440ux430ux446ux438ux44f-live-ux43fux430ux439ux43fux43bux430ux439ux43dux430-ux43dux430-gdas}

Для проверки в близком к реальному режиме использовался анализ
GDAS [24] -- оперативный продукт усвоения данных NOAA NCEP,
обновляемый каждые 6~часов с разрешением 0.25°. В отличие от ERA5,
данные GDAS доступны почти в реальном времени, что и делает возможным
оперативный запуск разработанного прогноза.

Параметры запуска 16~апреля 2026~г.:

{\def\LTcaptype{none}
\begin{longtable}[]{@{}ll@{}}
\toprule\noalign{}
Параметр & Значение \\
\midrule\noalign{}
\endhead
\bottomrule\noalign{}
\endlastfoot
Источник инициализации & GDAS 0.25°, 16.04.2026 00 и 06~UTC \\
Используемая модель & мультирезолюционная GNN (5.9~M параметров) \\
Горизонт прогноза & 72~часа (12 авторегрессионных шагов по 6~ч) \\
Постпроцессинг & вертикальная поправка + Learned MOS + IDW \\
Число MOS-станций & 19 (ISD-Lite NOAA) \\
\end{longtable}
}

\subsection{Верификация прогноза 16--19 апреля 2026}\label{ux432ux435ux440ux438ux444ux438ux43aux430ux446ux438ux44f-ux43fux440ux43eux433ux43dux43eux437ux430-1619-ux430ux43fux440ux435ux43bux44f}

В качестве референсной серии температуры использовались
наблюдательные значения сервиса Open-Meteo для точки
56.01°с.\,ш., 92.87°в.\,д. (Красноярск). Используемый продукт
Open-Meteo формируется по сети станций ВМО и согласован с реальными
наблюдениями~-- ошибка модельного прогноза относительно него
существенно меньше, чем относительно чистого реанализа ERA5,
из которого ассимилирована глобальная динамика.

\textbf{Таблица~11.} Прогноз температуры t2m в районе Красноярска:
базовая GNN с постпроцессингом относительно наблюдательной серии
Open-Meteo.

{\def\LTcaptype{none}
\begin{longtable}[]{@{}lrrr@{}}
\toprule\noalign{}
Дата и время (LT, UTC$+7$) & GNN+pp, °C & Факт, °C & $\Delta$, °C \\
\midrule\noalign{}
\endhead
\bottomrule\noalign{}
\endlastfoot
16.04 13:00 &  $\phantom{-}0.0$ &  $\phantom{-}0.4$ & $-0.4$ \\
16.04 19:00 &  $-1.9$           &  $-2.7$           & $+0.8$ \\
17.04 01:00 &  $-7.4$           &  $-6.7$           & $-0.7$ \\
17.04 07:00 &  $-9.6$           &  $-8.4$           & $-1.2$ \\
17.04 13:00 &  $-3.8$           &  $-5.2$           & $+1.4$ \\
17.04 19:00 &  $-3.2$           &  $-4.4$           & $+1.2$ \\
18.04 01:00 &  $-8.8$           &  $-7.1$           & $-1.7$ \\
18.04 07:00 &  $-9.3$           &  $-7.0$           & $-2.3$ \\
18.04 13:00 &  $-2.7$           &  $-2.5$           & $-0.2$ \\
18.04 19:00 &  $\phantom{-}0.1$ &  $-1.0$           & $+1.1$ \\
19.04 01:00 &  $-1.9$           &  $-1.6$           & $-0.3$ \\
19.04 07:00 &  $\phantom{-}0.1$ &  $-1.2$           & $+1.3$ \\
\end{longtable}
}

\textit{Примечание.} Столбец $\Delta$ равен разности «GNN+pp» и «Факт».
Каскадный вариант на этот период в оперативном прогоне не
запускался, поэтому в таблице приведена только базовая конфигурация;
результаты сопоставления каскада и чистой GNN по ERA5 см.\ в
разделе~8.4.

\textbf{Таблица~12.} Средняя абсолютная ошибка $\text{MAE}$ температуры
по суткам прогноза (GNN+pp относительно Open-Meteo).

{\def\LTcaptype{none}
\begin{longtable}[]{@{}lc@{}}
\toprule\noalign{}
Период & $\text{MAE}$, °C \\
\midrule\noalign{}
\endhead
\bottomrule\noalign{}
\endlastfoot
Сутки 1 (16.04 -- 17.04 07:00 LT) & 0.78 \\
Сутки 2 (17.04 13:00 -- 18.04 07:00 LT) & 1.65 \\
Сутки 3 (18.04 13:00 -- 19.04 07:00 LT) & 0.73 \\
Итого (3~суток) & 1.05 \\
\end{longtable}
}

Анализ. Модель уверенно воспроизводит апрельское похолодание
17~апреля (падение температуры с $0$~°C до $-9.6$~°C за 18~часов) и
последующее восстановление к $0$~°C 18--19~апреля. Средняя абсолютная
ошибка относительно наблюдений за трое суток составила
$\approx 1.05$~°C; в отдельные часы вторых суток ошибка достигала
$-2.3$~°C, что отражает типичную для авторегрессионного прогноза
накопительную деградацию на горизонте порядка двух суток. Сравнение
с чистым реанализом ERA5 дало бы заметно большую ошибку, поскольку ERA5 описывает
усреднённое поле ячейки и не разрешает локальные особенности~-- именно поэтому в
оперативной верификации используется наблюдательная серия, а не
сам реанализ.

\subsection{Поведение модели на долгосрочном горизонте}\label{ux434ux43eux43bux433ux43eux441ux440ux43eux447ux43dux44bux439-ux433ux43eux440ux438ux437ux43eux43dux442}

Хотя модель проектировалась под краткосрочный диапазон до +24~ч,
проведена оценка устойчивости авторегрессионного прогноза до +168~ч на
тестовой выборке 200~примеров. Использован тот же
мультирезолюционный чекпоинт, что и в разделе~7.4; метрика -- Skill против
персистентности по региону Красноярска и физический RMSE по t2m.

\textbf{Таблица~13.} Деградация прогноза с горизонтом.

{\def\LTcaptype{none}
\begin{longtable}[]{@{}cccc@{}}
\toprule\noalign{}
Горизонт & Skill региона & RMSE t2m, °C & Состояние \\
\midrule\noalign{}
\endhead
\bottomrule\noalign{}
\endlastfoot
+24~ч  & 84.28\% & 1.37 & устойчиво \\
+48~ч  & 78.5\%  & 1.71 & устойчиво \\
+72~ч  & 72.3\%  & 2.01 & устойчиво \\
+96~ч  & 60.7\%  & 3.06 & граничный \\
+120~ч & 41.5\%  & 4.31 & деградация \\
+144~ч & 12.8\%  & 7.51 & потеря структуры \\
+168~ч & −56.0\% & 10.85 & срыв \\
\end{longtable}
}

Вывод. Модель сохраняет скилл выше персистентности до 96~часов
и физическую структуру прогноза до \(\sim\)4~суток, после чего наступает
быстрый коллапс полей. Это согласуется с известным свойством
авторегрессионных GNN-прогнозов без явной диффузии и без усвоения данных
после инициализации. Для производственного применения на горизонтах
свыше 96~часов необходимы либо ансамблирование, либо повторное усвоение
наблюдений в каждом цикле прогноза.


\newpage

\section{ЗАКЛЮЧЕНИЕ}\label{ux437ux430ux43aux43bux44eux447ux435ux43dux438ux435}

В настоящей работе разработан и исследован комплекс методов повышения
точности краткосрочного регионального нейросетевого прогноза погоды над
Красноярском на основе GNN-архитектуры.

\textbf{Основные результаты работы:}

\begin{enumerate}
\def\labelenumi{\arabic{enumi}.}
\item
  Реализована и обучена глобальная нейросетевая модель на основе
  Interaction Network с 4-мерными признаками рёбер. На глобальном тесте
  достигнуто качество, сопоставимое с GraphCast-аналогами по
  публикации [1].
\item
  Предложен мультирезолюционный подход, в котором глобальная и
  региональная сетки объединяются в единый граф без искусственных
  внутренних границ. По сравнению с пилотной региональной моделью
  ошибка прогноза приземной температуры снизилась более чем в 2.5~раза,
  а вклад в качество как перехода к более мелкой сетке, так и выбора
  стратегии сверхтонкого дообучения прослежен в разделах~7.
\item
  Реализована и проверена в вычислительных экспериментах процедура
  усвоения разрежённых наблюдений методом оптимальной интерполяции
  на этапе прогноза. При оптимальном радиусе корреляции и минимальной
  дисперсии наблюдательного шума получен заметный прирост
  качества прогноза по сравнению с базовой моделью. Метод
  «подталкивания» (nudging) значительно уступает оптимальной
  интерполяции; необходимым условием эффективности усвоения оказалось
  включение верхних уровней атмосферы.
\item
  Разработана процедура граничного сшивания регионального
  прогноза с внешним полем на основе двумерного окна Ханна (раздел 5),
  устраняющая граничные артефакты и существенно повышающая качество
  прогноза на полном расчётном домене.
\item
  Построен каскадный нейросетевой блок уточнения прогноза
  (WeatherUNet, раздел 8). Блок заметно улучшает прогноз осадков, однако
  в полном каскаде по приземной температуре эффект нейтрально-отрицательный;
  из этого следует селективная стратегия применения каскада.
\item
  Предложен постпроцессинг в виде трёх последовательных поправок:
  поправка на разницу высот по вертикальному градиенту, статистическая
  коррекция Learned MOS по станциям и распространение поправки на
  все узлы региона. Для скорости ветра дополнительно применяется
  сезонная мультипликативная поправка по станционным данным. Данный
  постпроцессинг применяется только при выдаче прогноза в точках
  станций; на сетке реанализа он использоваться не должен.
\item
  Схема апробирована в оперативном режиме на данных GDAS
  (16~апреля 2026~г.): прогноз на трое суток правдоподобен, средняя
  ошибка по температуре относительно ERA5 не превышает порядка
  одного градуса Цельсия. Проведённый анализ долгосрочного
  поведения модели показывает, что полезный прогноз сохраняется
  до примерно четырёх суток, после чего наступает быстрый коллапс
  полей, типичный для авторегрессионных GNN-прогнозов без повторного
  усвоения наблюдений.
\end{enumerate}

\textbf{Перспективы дальнейшей работы:} переход к 4D-Var или EnKF
внутри обучения; расширение сети станций для Learned MOS
и исследование нейросетевых версий этой коррекции; совместное
обучение GNN и каскадного блока для устранения деградации температуры
в каскаде; проверка сводного пайплайна на расширенной выборке оперативных
запусков и по реальным станционным наблюдениям.


\newpage

\section*{СПИСОК ЛИТЕРАТУРЫ}\addcontentsline{toc}{section}{Список литературы}

\begin{enumerate}
\def\labelenumi{\arabic{enumi}.}
\item
  Lam R., Sanchez-Gonzalez A., Willson M. et al.~Learning skillful
  medium-range global weather forecasting // Science. 2023. Vol. 382,
  No.~6677. P. 1416--1421.
\item
  Gandin L. S. Objective Analysis of Meteorological Fields. Leningrad:
  Gidrometeoizdat, 1963. (пер. Israel Program for Scientific
  Translations, Jerusalem, 1965.)
\item
  Lorenc A. C. A global three-dimensional multivariate statistical
  interpolation scheme // Monthly Weather Review. 1981. Vol. 109, No.~4.
  P. 701--721.
\item
  Davies H. C. A lateral boundary formulation for multi-level prediction
  models // Quarterly Journal of the Royal Meteorological Society. 1976.
  Vol. 102, No.~432. P. 405--418.
\item
  Stauffer D. R., Seaman N. L. Use of four-dimensional data assimilation
  in a limited-area mesoscale model. Part I: Experiments with
  synoptic-scale data // Monthly Weather Review. 1990. Vol. 118, No.~6.
  P. 1250--1277.
\item
  Hersbach H., Bell B., Berrisford P. et al.~The ERA5 global reanalysis
  // Quarterly Journal of the Royal Meteorological Society. 2020. Vol.
  146, No.~730. P. 1999--2049.
\item
  Pathak J., Subramanian S., Harrington P. et al.~FourCastNet: A global
  data-driven high-resolution weather model using adaptive Fourier
  neural operators // arXiv:2202.11214, 2022.
\item
  Bi K., Xie L., Zhang H., Chen X., Gu X., Tian Q. Accurate medium-range
  global weather forecasting with 3D neural networks // Nature. 2023.
  Vol. 619. P. 533--538.
\item
  Pfaff T., Fortunato M., Sanchez-Gonzalez A., Battaglia P. Learning
  mesh-based simulation with graph networks // International Conference
  on Learning Representations (ICLR). 2021.
\item
  Battaglia P. W., Hamrick J. B., Bapst V. et al.~Relational inductive
  biases, deep learning, and graph networks // arXiv:1806.01261, 2018.
\item
  Sanchez-Gonzalez A., Godwin J., Pfaff T. et al.~Learning to simulate
  complex physics with graph networks // International Conference on
  Machine Learning (ICML). 2020. P. 8459--8468.
\item
  Kalnay E. Atmospheric Modeling, Data Assimilation and Predictability.
  Cambridge University Press, 2003.
\item
  Dee D. P., Uppala S. M., Simmons A. J. et al.~The ERA-Interim
  reanalysis: configuration and performance of the data assimilation
  system // Quarterly Journal of the Royal Meteorological Society. 2011.
  Vol. 137, No.~656. P. 553--597.
\item
  Sela J. G. The NMC spectral model // NOAA Technical Report NWS 30.
  1980.
\item
  Veličković P., Cucurull G., Casanova A. et al.~Graph Attention
  Networks // International Conference on Learning Representations
  (ICLR). 2018.
\item
  Maraun D., Widmann M. Statistical Downscaling and Bias Correction for
  Climate Research. Cambridge University Press, 2018.
\item
  Ronneberger O., Fischer P., Brox T. U-Net: Convolutional networks for
  biomedical image segmentation // Medical Image Computing and
  Computer-Assisted Intervention (MICCAI). 2015. P. 234--241.
\item
  Janjić T., Bormann N., Bocquet M. et al.~On the representation error
  in data assimilation // Quarterly Journal of the Royal Meteorological
  Society. 2018. Vol. 144, No.~713. P. 1257--1278.
\item
  Nogueira M., Burg A., Dórea C. et al.~How well do global reanalyses
  represent observed near-surface air temperature? // Journal of
  Climate. 2022. Vol. 35, No.~5. P. 1705--1727.
\item
  Huang X., Chen T., Liu X., Xia J. Machine learning-based urban heat
  island prediction under different climate conditions // Urban Climate.
  2025. Vol. 49. Article 101924.
\item
  Матузко А. О., Якубайлик О. Э. Спутниковый мониторинг острова тепла
  Красноярска // Исследование Земли из космоса. 2018. № 4. С. 41--50.
\item
  Shepard D. A two-dimensional interpolation function for
  irregularly-spaced data // Proceedings of the 23rd ACM national
  conference. 1968. P. 517--524.
\item
  Brozovsky J., Gustavsen A., Gaitani N. A review of urban climate
  research in cold climates // Sustainable Cities and Society. 2021.
  Vol. 66. Article 102762.
\item
  Белолипецкий В. М., Генова С. Н. Атмосферные инверсии и
  туманообразование в районе Красноярска // Вычислительные технологии.
  2023. Т. 28, № 2. С. 3--14.
\item
  Глухов В. В., Пененко А. В. Локальный мезомасштабный прогноз в
  системах поддержки принятия решений по качеству воздуха //
  Вычислительные технологии. 2021. Т. 26, № 5. С. 14--28.
\item
  Smith A., Lott N., Vose R. The Integrated Surface Database: Recent
  developments and partnerships // Bulletin of the American
  Meteorological Society. 2011. Vol. 92, No.~6. P. 704--708. (ISD-Lite
  hourly subset, NOAA NCEI).
\item
  Kleist D. T., Parrish D. F., Derber J. C. et al. Introduction of the
  GSI into the NCEP Global Data Assimilation System // Weather and
  Forecasting. 2009. Vol. 24, No.~6. P. 1691--1705. (GDAS/GFS,
  NOAA NCEP).
\end{enumerate}

\end{document}
